%==============================================================================
%  TRABAJO DE GRADO - Schedule de ejecución basado en portafolios para el
%  problema del Job Shop clásico
%  Universidad del Valle - Escuela de Ingeniería de Sistemas y Computación
%
%  Estructura:
%    main.tex                -> preámbulo, portada y orden del documento
%    preliminares/          -> agradecimientos y resumen
%    capitulos/             -> un archivo por capítulo
%    anexos/                -> un archivo por anexo
%    referencias.tex        -> lista de referencias (formato APA)
%
%  Compilar con pdfLaTeX (Overleaf: Menu -> Compiler -> pdfLaTeX).
%==============================================================================
\documentclass[12pt, letterpaper, oneside]{report}

%------------------------------------------------------------------------------
%  Idioma y codificación
%------------------------------------------------------------------------------
\usepackage[utf8]{inputenc}
\usepackage[T1]{fontenc}
\usepackage[spanish, es-tabla, es-nodecimaldot]{babel}

%------------------------------------------------------------------------------
%  Tipografía, márgenes e interlineado
%------------------------------------------------------------------------------
\usepackage{mathptmx}                 % Times
\usepackage[margin=2.54cm]{geometry}
\usepackage{setspace}
\onehalfspacing
\setlength{\parindent}{1.27cm}
\usepackage[expansion=false]{microtype}

%------------------------------------------------------------------------------
%  Paquetes generales
%------------------------------------------------------------------------------
\usepackage{amsmath, amssymb}
\usepackage{graphicx}
\graphicspath{{figuras/}}
\usepackage{booktabs}
\usepackage{array}
\usepackage{tabularx}
\usepackage{longtable}
\usepackage{enumitem}
\setlist{nosep}
\usepackage{caption}
\usepackage{float}               % [H]: cada tabla y figura queda donde se escribe, sin cortar párrafos
\captionsetup{font=small, labelfont=bf, labelsep=period}
% Las páginas que solo contienen tablas o figuras se alinean arriba.
\makeatletter
\setlength{\@fptop}{0pt}
\makeatother

%------------------------------------------------------------------------------
%  Citas en formato autor-año (APA)
%------------------------------------------------------------------------------
\usepackage[round, authoryear]{natbib}
\setlength{\bibhang}{1.27cm}
\setlength{\bibsep}{0.5em}

%------------------------------------------------------------------------------
%  Encabezados de capítulo y sección
%------------------------------------------------------------------------------
\usepackage{titlesec}
\titleformat{\chapter}[display]
  {\normalfont\bfseries\centering}{\chaptertitlename\ \thechapter}{0.4em}{}
\titleformat{name=\chapter, numberless}
  {\normalfont\bfseries\centering}{}{0pt}{}
\titlespacing*{\chapter}{0pt}{0pt}{1.5em}
\titleformat{\section}
  {\normalfont\bfseries\centering}{\thesection}{0.6em}{}
\titlespacing*{\section}{0pt}{1.2em}{0.6em}
\titleformat{\subsection}
  {\normalfont\bfseries}{\thesubsection}{0.6em}{}
\titlespacing*{\subsection}{0pt}{1em}{0.4em}
\titleformat{\subsubsection}
  {\normalfont\bfseries\itshape}{\thesubsubsection}{0.6em}{}

%------------------------------------------------------------------------------
%  Numeración de página en la esquina superior derecha
%------------------------------------------------------------------------------
\usepackage{fancyhdr}
\fancypagestyle{plain}{%
  \fancyhf{}%
  \fancyhead[R]{\textsc{\thepage}}%
  \renewcommand{\headrulewidth}{0pt}%
}
\pagestyle{plain}
\setlength{\headheight}{14.5pt}

%------------------------------------------------------------------------------
%  Nombres de los elementos del documento
%------------------------------------------------------------------------------
\addto\captionsspanish{%
  \renewcommand{\contentsname}{Índice general}%
  \renewcommand{\listfigurename}{Índice de figuras}%
  \renewcommand{\listtablename}{Índice de tablas}%
  \renewcommand{\bibname}{Referencias}%
  \renewcommand{\appendixname}{Anexo}%
}

%------------------------------------------------------------------------------
%  Enlaces (al final del preámbulo)
%------------------------------------------------------------------------------
\usepackage[hidelinks]{hyperref}
\urlstyle{same}

%------------------------------------------------------------------------------
%  Datos del trabajo
%------------------------------------------------------------------------------
\newcommand{\titulo}{SCHEDULE DE EJECUCIÓN BASADO EN PORTAFOLIOS PARA EL\\
  PROBLEMA DEL JOB SHOP CLÁSICO}
\newcommand{\fecha}{Noviembre de 2026}

%==============================================================================
\begin{document}
%==============================================================================

%------------------------------------------------------------------------------
%  Portada
%------------------------------------------------------------------------------
\pagenumbering{Alph}
\begin{titlepage}
  \centering
  \singlespacing
  \includegraphics[width=4cm]{logo_univalle.png}\par
  \vspace{1cm}
  {\large\bfseries \titulo\par}
  \vspace{0.8cm}
  Modalidad Trabajo Profesional\par
  \vfill
  Oscar David Imbachi Ordoñez\\
  201931037\\
  \texttt{oscar.imbachi@correounivalle.edu.co}\par
  \vfill
  Director\\
  Robinson Andrey Duque Agudelo, Ph.D.\\
  Universidad del Valle\\
  \texttt{robinson.duque@correounivalle.edu.co}\par
  \vspace{0.8cm}
  Codirector\\
  Alejandro Arbeláez Rodríguez, Ph.D.\\
  Universitat Politècnica de València\\
  \texttt{aarbrod@upv.es}\par
  \vfill
  Universidad del Valle\\
  Facultad de Ingeniería\\
  Escuela de Ingeniería de Sistemas y Computación\\
  Programa Académico de Ingeniería de Sistemas (3743)\\
  Santiago de Cali\\
  \fecha\par
\end{titlepage}

%------------------------------------------------------------------------------
%  Páginas preliminares (numeración romana)
%------------------------------------------------------------------------------
\pagenumbering{roman}
\input{preliminares/agradecimientos}
\input{preliminares/resumen}

\tableofcontents
\listoftables
\listoffigures

%------------------------------------------------------------------------------
%  Cuerpo del documento (numeración arábiga)
%------------------------------------------------------------------------------
\cleardoublepage
\pagenumbering{arabic}
%==============================================================================
\chapter{Introducción}
\label{cap:introduccion}
%==============================================================================

%------------------------------------------------------------------------------
\section{Contexto y planteamiento del problema}
\label{sec:contexto}
%------------------------------------------------------------------------------
% Fuente: Anteproyecto, Cap. 1 (Introducción) y Cap. 2 (Planteamiento).

El Job Shop Scheduling Problem (JSSP) es uno de los problemas más estudiados en
el campo de la optimización combinatoria. En este problema, un conjunto de
trabajos debe ser procesado en distintas máquinas, cumpliendo para cada trabajo
una secuencia específica de operaciones. El objetivo suele ser minimizar el
tiempo total de finalización (\textit{makespan}), aunque pueden considerarse
otros criterios como el retraso promedio o la utilización de recursos
\citep{pinedo2022}.

El JSSP es un problema NP-duro, por lo que su resolución eficiente requiere
algoritmos avanzados, entre ellos heurísticas, metaheurísticas y métodos
híbridos. Sin embargo, el rendimiento de estos algoritmos depende en gran
medida de las características de cada instancia, lo que hace difícil determinar
de antemano cuál ofrecerá los mejores resultados. En la práctica, seleccionar
manualmente el algoritmo más apropiado es ineficiente y requiere experiencia
especializada.

Para abordar esta incertidumbre, los portafolios de algoritmos se han
convertido en una estrategia efectiva, puesto que combinan múltiples métodos y
seleccionan aquel que se espera tenga el mejor desempeño
\citep{smithmiles2009}. No obstante, la mayoría de los enfoques adoptan la
política \textit{Winner Takes All}, que consiste en seleccionar un solo
algoritmo por instancia. Este enfoque es vulnerable a errores de selección y no
aprovecha la complementariedad entre algoritmos. Por ello, junto a selectores por
instancia como SATzilla \citep{xu2008satzilla}, han surgido métodos como
CPHydra \citep{omahony2008cphydra} y SUNNY \citep{amadini2014sunny}, que
emplean técnicas de aprendizaje automático para definir un \textit{schedule}
de ejecución, es decir, el orden y el tiempo asignado a cada algoritmo dentro
del portafolio. Estos sistemas han demostrado
mejorar el rendimiento global y la robustez en la resolución de problemas
complejos.

Este trabajo aplica esa idea al JSSP clásico mediante el diseño de un
portafolio de algoritmos y un esquema de ejecución planificado con
\textit{Machine Learning}, capaz de decidir dinámicamente la secuencia y la
duración de ejecución de los métodos más apropiados para cada instancia.

%------------------------------------------------------------------------------
\section{Formulación del problema}
\label{sec:formulacion}
%------------------------------------------------------------------------------
% Fuente: Anteproyecto, Cap. 2 (Pregunta problema).

La pregunta que guía el trabajo es la siguiente:

\begin{quote}
  ¿Cómo diseñar un \textit{schedule} de ejecución basado en portafolios que
  determine el orden y tiempo óptimo de ejecución de algoritmos, utilizando
  \textit{Machine Learning}, para mejorar la eficiencia y robustez en la
  resolución del Job Shop Scheduling Problem clásico?
\end{quote}

Para responderla, el trabajo construye un portafolio de algoritmos
complementado por un modelo de selección automática que indica en qué orden
ejecutar los algoritmos y durante cuánto tiempo, es decir, el \textit{schedule}
de ejecución.

%------------------------------------------------------------------------------
\section{Justificación}
\label{sec:justificacion}
%------------------------------------------------------------------------------
% Fuente: Anteproyecto, Cap. 3.1 (económica) y 3.2 (académica).

En el ámbito académico, el trabajo aporta al estudio del problema de selección
de algoritmos (\textit{Algorithm Selection Problem}), una línea de
investigación que combina el aprendizaje automático y la optimización
combinatoria. Ofrece un entorno experimental para explorar la caracterización
de instancias del JSSP, la aplicación de técnicas de clasificación y regresión
para predecir el rendimiento de los algoritmos, y la integración de modelos de
aprendizaje con métodos de optimización tradicionales, como la búsqueda tabú,
los algoritmos genéticos o el recocido simulado.

Desde la perspectiva económica y productiva, una planificación ineficiente en
entornos industriales genera costos por tiempos muertos, sobreutilización de
recursos y retrasos en la entrega de productos. Un portafolio inteligente puede
reducir los costos operativos mediante una mejor asignación de recursos,
disminuir el tiempo total de procesamiento (\textit{makespan}) e incrementar la
productividad de la planta. En el plano computacional, seleccionar
adecuadamente el algoritmo reduce tanto el tiempo de cómputo como el consumo de
recursos.

%------------------------------------------------------------------------------
\section{Objetivos}
\label{sec:objetivos}
%------------------------------------------------------------------------------
% Fuente: Anteproyecto, Cap. 4.1 y 4.2 (texto aprobado, no modificar sin
% autorización del director).

\subsection{Objetivo general}

Desarrollar un \textit{schedule} de ejecución basado en portafolios para el Job
Shop Scheduling Problem clásico, utilizando técnicas de \textit{Machine
Learning}.

\subsection{Objetivos específicos}

\begin{enumerate}[label=\arabic*., leftmargin=1.27cm, itemsep=0.3em]
  \item Analizar y comparar distintos enfoques de portafolios de algoritmos
    aplicados al Job Shop Scheduling Problem, incluyendo estrategias
    \textit{Winner Takes All} y planificación secuencial.
  \item Identificar y extraer características relevantes de las instancias del
    JSSP.
  \item Implementar un portafolio que integre la predicción y selección
    automática del algoritmo más adecuado.
  \item Validar experimentalmente el portafolio mediante pruebas con instancias
    de referencia.
\end{enumerate}

%------------------------------------------------------------------------------
\section{Alcance y limitaciones}
\label{sec:alcance}
%------------------------------------------------------------------------------
% Fuente: Anteproyecto, Cap. 5.

El trabajo se centra en el estudio y la resolución del Job Shop Scheduling
Problem clásico, con la minimización del \textit{makespan} como criterio
principal de optimización. Para ello se trabaja con un conjunto limitado pero
representativo de algoritmos, que abarca métodos heurísticos, metaheurísticos y
exactos, con el propósito de evaluar su comportamiento ante diferentes tipos de
instancias. El enfoque se basa en modelos de aprendizaje supervisado, que se
entrenan y validan con conjuntos de datos de referencia.

El portafolio se limita a la versión clásica del problema. No incluye variantes
como el Flexible Job Shop ni las versiones multiobjetivo, que se consideran
líneas de trabajo futuro.

En la práctica, el conjunto de algoritmos quedó compuesto por dos
\textit{solvers} exactos de programación por restricciones, Chuffed y Gecode,
en tres configuraciones, junto con una heurística de despacho que calcula una
cota superior del \textit{makespan} y una solución de respaldo. OR-Tools CP-SAT
se usa como \textit{solver} de referencia fuera del portafolio. La evaluación se
hace sobre 242 instancias clásicas de la literatura con un límite de 300
segundos por instancia. Las limitaciones que se derivan de estas decisiones se
discuten en la Sección~\ref{sec:limitaciones}.

%------------------------------------------------------------------------------
\section{Resultados obtenidos y contribuciones}
\label{sec:resultados_contribuciones}
%------------------------------------------------------------------------------
% Fuente: Anteproyecto, Cap. 4.3 y Tabla 4.1. Al cierre del trabajo, la
% columna de resultados esperados se reemplaza por los resultados obtenidos.

La Tabla~\ref{tab:objetivos_resultados} relaciona cada objetivo específico con
el resultado obtenido y con el capítulo de este documento donde se presenta.

\begin{table}[H]
  \centering
  \small
  \singlespacing
  \caption{Objetivos específicos, resultados obtenidos y capítulo donde se presentan.}
  \label{tab:objetivos_resultados}
  \begin{tabularx}{\textwidth}{>{\raggedright\arraybackslash}p{0.30\textwidth}
                               >{\raggedright\arraybackslash}X
                               >{\centering\arraybackslash}p{0.13\textwidth}}
    \toprule
    \textbf{Objetivo específico} & \textbf{Resultado obtenido} & \textbf{Capítulo} \\
    \midrule
    1. Analizar y comparar enfoques de portafolios (\textit{Winner Takes All} y
    planificación secuencial)
      & Comparación de \textit{Winner Takes All}, SUNNY, un \textit{schedule}
        estático y un \textit{schedule} con aprendizaje automático, en calidad
        de solución, óptimos probados, tiempo (PAR10) y robustez.
      & \ref{cap:portafolio}, \ref{cap:resultados} \\
    \addlinespace
    2. Identificar y extraer características de las instancias del JSSP
      & 65 características de \texttt{mzn2feat} y 20 propias del JSSP, junto
        con una base de conocimiento que incluye el desempeño y las
        trayectorias de cada configuración.
      & \ref{cap:caracterizacion} \\
    \addlinespace
    3. Implementar un portafolio con predicción y selección automática del
    algoritmo
      & Prototipo que calcula cotas, extrae características y construye el
        orden y el tiempo de ejecución por instancia con tres modelos de
        aprendizaje automático.
      & \ref{cap:portafolio} \\
    \addlinespace
    4. Validar experimentalmente el portafolio con instancias de referencia
      & Validación \textit{leave-one-out} sobre 236 instancias y tabla
        comparativa entre las configuraciones individuales con y sin cotas,
        los métodos de portafolio y el VBS.
      & \ref{cap:diseno}, \ref{cap:resultados} \\
    \bottomrule
  \end{tabularx}
\end{table}

Las principales contribuciones del trabajo son:

\begin{itemize}
  \item Un modelo con cotas para el JSSP en MiniZinc, que combina una cota
    inferior con la solución de la heurística de Giffler y Thompson y reduce el
    \textit{gap} medio de Chuffed de 648\,\% a 17\,\%. La comparación con los
    mismos \textit{solvers} sin cotas muestra que, sin ellas, ni siquiera
    eligiendo en cada instancia la mejor configuración el \textit{gap} medio
    baja del 24.5\,\%, mientras que con cotas todos los enfoques de portafolio
    quedan cerca del 13\,\%.
  \item Una base de conocimiento con trayectorias, que registra cuándo encontró
    cada configuración cada solución y permite evaluar cualquier
    \textit{schedule} sin volver a ejecutar los \textit{solvers}.
  \item Un conjunto de 20 características propias del JSSP, que se calculan en
    menos de un segundo y obtienen resultados comparables a los de
    \texttt{mzn2feat}.
  \item Una comparación controlada de estrategias de portafolio, que muestra
    que, con este portafolio y estas instancias, la selección de una
    configuración es mejor en
    promedio y el \textit{schedule} es más robusto frente a errores de
    selección.
  \item Un repositorio reproducible con las instancias, los modelos, el código,
    los resultados de todas las campañas y los programas que generan las tablas
    y figuras de este documento.
\end{itemize}

%------------------------------------------------------------------------------
\section{Organización del documento}
\label{sec:organizacion}
%------------------------------------------------------------------------------

El resto del documento se organiza en siete capítulos y cuatro anexos. El
Capítulo~\ref{cap:marco} presenta el marco de referencia sobre el JSSP, la
programación por restricciones, el problema de selección de algoritmos, los
portafolios de algoritmos, el uso del aprendizaje automático para predecir el
desempeño de los algoritmos y el estado del arte. El
Capítulo~\ref{cap:problema} define formalmente el problema y la metodología de
trabajo. El Capítulo~\ref{cap:caracterizacion} describe las instancias, la
extracción de características y la construcción del conjunto de datos. El
Capítulo~\ref{cap:portafolio} presenta el diseño y la implementación del
portafolio, incluidas las estrategias \textit{Winner Takes All} y de
\textit{schedule} secuencial. El Capítulo~\ref{cap:diseno} expone el diseño
experimental y el Capítulo~\ref{cap:resultados}, los resultados y su análisis.
Finalmente, el Capítulo~\ref{cap:conclusiones} recoge las conclusiones, las
limitaciones y el trabajo futuro. Los anexos contienen los resultados por
instancia, el catálogo de características, la configuración de los
experimentos y una guía de reproducción.

%==============================================================================
\chapter{Marco de referencia}
\label{cap:marco}
%==============================================================================

La construcción de un portafolio para el JSSP articula conceptos de la
programación de operaciones en planta, la programación por restricciones y la
selección automática de algoritmos mediante aprendizaje supervisado. A
continuación se revisan estos tres fundamentos y la revisión sistemática de
literatura que delimita la brecha abordada.

%------------------------------------------------------------------------------
\section{El Job Shop Scheduling Problem}
\label{sec:jssp}
%------------------------------------------------------------------------------
% Fuente: Anteproyecto, Cap. 6.1.1.

El Job Shop Scheduling Problem (JSSP) es uno de los problemas más estudiados
de la programación de tareas en entornos productivos. Consiste en asignar un
conjunto de $n$ trabajos, cada uno compuesto por una secuencia de
operaciones, a un conjunto de $m$ máquinas, de manera que cada operación se
procese en una máquina específica siguiendo un orden predefinido. El objetivo
clásico es minimizar el tiempo total de finalización (\textit{makespan})
\citep{pinedo2022}.

El problema es NP-duro, puesto que \citet{garey1976} demostraron que su versión
de decisión es NP-completa. Como cada máquina puede procesar sus $n$
operaciones en $n!$ órdenes distintos, el número de secuencias posibles crece
como $(n!)^m$, y por eso instancias de apenas $10 \times 10$, como
\texttt{ft10}, permanecieron sin resolver durante más de dos décadas
\citep{fisher1963, carlier1989}. En la práctica, los métodos exactos solo
garantizan el óptimo en instancias pequeñas o medianas dentro de un tiempo
razonable, y el desempeño de un mismo método varía mucho de una instancia a
otra, lo que motiva el uso de portafolios (Sección~\ref{sec:portafolios}).

Los métodos de solución suelen agruparse en tres familias
\citep{pinedo2022, strassl2022}. Los \textbf{métodos exactos} garantizan el óptimo si se
les da tiempo suficiente, e incluyen la ramificación y acotamiento específica
para el problema \citep{brucker1994}, la programación lineal entera mixta y la
programación por restricciones, que ha mostrado ser especialmente competitiva
en el JSSP \citep{ku2016}. Las \textbf{heurísticas constructivas} producen una
solución en un solo paso aplicando reglas de prioridad, como atender primero
la operación más corta \citep{panwalkar1977}. Entre ellas se encuentran el
algoritmo de \citet{giffler1960}, que genera programas activos, y el cuello de
botella móvil (\textit{shifting bottleneck}) de \citet{adams1988}. Son rápidas,
pero no garantizan la calidad de la solución. Las \textbf{metaheurísticas},
como la búsqueda tabú \citep{nowicki1996} o los algoritmos genéticos
\citep{yamada1992}, exploran el espacio de soluciones con estrategias
generales y suelen dar soluciones de muy buena calidad en instancias grandes,
aunque no prueban optimalidad. El portafolio de este trabajo se compone de
métodos exactos de programación por restricciones, y la heurística de
Giffler y Thompson se usa para acotar el problema
(Sección~\ref{sec:cotas}).

La investigación en el JSSP se apoya en instancias de referencia publicadas a
lo largo de varias décadas, entre las que destacan las de \citet{fisher1963},
\citet{lawrence1984}, \citet{adams1988}, \citet{applegate1991},
\citet{storer1992}, \citet{yamada1992}, \citet{taillard1993} y
\citet{demirkol1998}. Muchas se distribuyen a través de la OR-Library
\citep{beasley1990}, y sus cotas y mejores soluciones conocidas se actualizan
a medida que se publican nuevos resultados \citep{weise2019jssp}. Son las ocho
familias que se usan en este trabajo (Sección~\ref{sec:instancias}).

%------------------------------------------------------------------------------
\section{Programación por restricciones}
\label{sec:cp}
%------------------------------------------------------------------------------

En la programación por restricciones, el problema se describe de forma
declarativa mediante variables, dominios y restricciones, y un
\textit{solver} genérico se encarga de encontrar las soluciones
\citep{rossi2006}. Un problema de satisfacción de restricciones (CSP) se
define por una tripleta $(X, D, C)$ compuesta por variables $X$, un dominio
$D(x)$ de valores posibles para cada variable y restricciones $C$ que limitan
las combinaciones permitidas. Un problema de optimización con restricciones
(COP) añade una función objetivo. El JSSP con minimización del
\textit{makespan} es un COP cuyas variables son los tiempos de inicio de las
operaciones (Sección~\ref{sec:notacion}).

MiniZinc es un lenguaje de modelado de alto nivel, independiente del
\textit{solver} \citep{nethercote2007}, que separa el modelo (archivo
\texttt{.mzn}) de los datos de cada instancia (archivo \texttt{.dzn}). Antes
de resolver, el compilador traduce ambos a FlatZinc, el lenguaje de bajo nivel
que entienden los \textit{solvers}, y en esa traducción cada uno puede usar su
propia implementación de las restricciones globales, como
\texttt{disjunctive}, que impide que dos tareas usen a la vez el mismo
recurso. Gracias a esta separación, un único modelo del JSSP sirve para
ejecutar \textit{solvers} con técnicas muy distintas, y las diferencias de
desempeño se deben al \textit{solver}. MiniZinc es además el formato del
MiniZinc Challenge, la competencia anual de \textit{solvers} de restricciones
\citep{stuckey2014}.

Los \textit{solvers} clásicos combinan \textbf{propagación} y
\textbf{búsqueda} \citep{rossi2006}. La propagación elimina de los dominios
los valores que no pueden formar parte de ninguna solución. Por ejemplo, si
una operación no puede empezar antes del instante 30, su sucesora en el mismo
trabajo tampoco puede empezar antes de que la primera termine. Cuando la
propagación ya no avanza, la búsqueda elige una variable, le asigna un valor
o divide su dominio y vuelve a propagar, retrocediendo cuando encuentra una
contradicción. En optimización, cada solución encontrada impone que la
siguiente sea mejor, hasta probar que no existe una mejor. Gecode
\citep{gecode2019} es un \textit{solver} de este tipo, con propagadores
eficientes para muchas restricciones globales, pero, a diferencia de los
\textit{solvers} con generación perezosa de cláusulas, no aprende de los
fallos \citep{chu2011}.

La generación perezosa de cláusulas (\textit{Lazy Clause Generation}, LCG)
combina la propagación con el aprendizaje de cláusulas de los
\textit{solvers} SAT \citep{ohrimenko2009}. Cada reducción de un dominio se
registra con una explicación en forma de cláusula. Cuando se produce un fallo,
el \textit{solver} deriva de esas explicaciones una cláusula que resume la
causa del conflicto y la añade al problema, de modo que el mismo fallo no se
repite en otras ramas. Chuffed \citep{chu2011} implementa esta idea, que es
especialmente eficaz en problemas de \textit{scheduling}
\citep{schutt2011}.

CP-SAT, el \textit{solver} de restricciones de OR-Tools \citep{perron2025},
combina en un mismo motor un núcleo SAT con generación perezosa de cláusulas,
relajaciones de programación lineal que aportan cotas al objetivo, técnicas
de presolución y búsqueda local, y por defecto ejecuta en paralelo varias
estrategias de búsqueda \citep{perron2023}, aunque en este trabajo se usó con un único hilo, en
igualdad de condiciones con los demás \textit{solvers}
(Sección~\ref{sec:ejecucion_solvers}). Por sus resultados en el MiniZinc
Challenge y en \textit{scheduling}, se usa aquí como referencia del estado
del arte.

%------------------------------------------------------------------------------
\section{El problema de selección de algoritmos}
\label{sec:asp}
%------------------------------------------------------------------------------
% Fuente: Anteproyecto, Cap. 6.1.3.

\citet{rice1976} planteó el problema de selección de algoritmos como el de
determinar qué algoritmo, entre un conjunto disponible, es el más adecuado
para una instancia específica. Su modelo distingue cuatro espacios: el de
problemas $\mathcal{P}$, que contiene las instancias; el de características
$\mathcal{F}$, donde cada instancia $i$ se representa por un vector $f(i)$; el
de algoritmos $\mathcal{A}$, y el de desempeño $\mathcal{Y}$, donde $y(a, i)$
mide qué tan bien resuelve $a$ la instancia $i$. Se trata de encontrar una
función $S: \mathcal{F} \rightarrow \mathcal{A}$ que elija el algoritmo con
mejor desempeño esperado, y en la práctica $S$ se aprende de instancias cuyo
desempeño con cada algoritmo ya se conoce \citep{kotthoff2014,
kerschke2019}. Iniciativas como COSEAL promueven metodologías y conjuntos de
datos comunes para comparar estrategias de selección \citep{coseal2024}, y
las competencias de selección de algoritmos han fijado métricas comunes
\citep{lindauer2019}. Este trabajo extiende ese planteamiento para que la
función produzca un \textit{schedule} de ejecución completo
(Sección~\ref{sec:formulacion_schedule}).

La calidad de un selector depende de que las características describan bien
lo que hace difícil una instancia. Las \textbf{características genéricas} se
calculan sobre la representación que recibe el \textit{solver}, sin conocer
el dominio. En programación por restricciones, \texttt{mzn2feat}
\citep{amadini2014mzn2feat} extrae características del modelo FlatZinc (95 en
su versión 1.0, \citealp{amadini2015sunnycp}), y CPHydra
\citep{omahony2008cphydra} usa un conjunto más reducido de características
del mismo tipo. Sirven para
cualquier problema, pero dependen de cómo esté escrito el modelo. Las
\textbf{características específicas} usan conocimiento del dominio, como
ocurre en el JSSP con la relación entre trabajos y máquinas, la variabilidad
de los tiempos de proceso, la carga de las máquinas y la estructura de las
rutas \citep{mirshekarian2016, strassl2022}. Este trabajo combina ambos tipos
(Sección~\ref{sec:features}).

%------------------------------------------------------------------------------
\section{Portafolios de algoritmos}
\label{sec:portafolios}
%------------------------------------------------------------------------------
% Fuente: Anteproyecto, Cap. 6.1.2.

Un portafolio de algoritmos es un conjunto de métodos para un mismo tipo de
problema, cada uno con ventajas y debilidades propias, del que se elige o
combina el más apropiado según la instancia \citep{gomes2001, kotthoff2014}. Su
motivación es que ningún método es el mejor en todos los casos, como lo
expresa el teorema \textit{No Free Lunch} \citep{wolpert1997}.

La estrategia más común, que aquí se llama \textit{Winner Takes All} (WTA)
como en el anteproyecto, elige un único algoritmo por instancia y le dedica
todo el tiempo disponible. El selector puede ser un clasificador que predice
el mejor algoritmo o un conjunto de modelos de regresión que predicen el
desempeño de cada uno \citep{xu2008satzilla, kotthoff2014}. Si el algoritmo
elegido falla, ya no queda tiempo para otro, y ese riesgo crece cuando las
predicciones son poco fiables.

Un \textit{schedule} secuencial de ejecución reparte el presupuesto $T$ entre
varios algoritmos que se ejecutan uno tras otro. A cambio de dar menos tiempo
a cada uno, reduce el riesgo de una mala selección, y suele bastar porque
muchas instancias se resuelven rápido con el algoritmo adecuado
\citep{omahony2008cphydra, amadini2014sunny}. Construirlo exige decidir qué
algoritmos incluir, cuánto tiempo darle a cada uno y en qué orden ejecutarlos,
que son las tres decisiones de la pregunta de investigación de este trabajo.

Entre los sistemas representativos, SATzilla \citep{xu2008satzilla} es un
portafolio WTA para SAT que predice el tiempo de cada \textit{solver} y elige
el menor, tras ejecutar brevemente unos \textit{presolvers}. Hydra
\citep{xu2010hydra} construye el portafolio configurando versiones de un
mismo algoritmo parametrizado que se complementen entre sí. CPHydra
\citep{omahony2008cphydra}, el primer portafolio de \textit{solvers} para
programación por restricciones \citep{amadini2015sunnycp}, busca las instancias más parecidas a la nueva
y calcula el \textit{schedule} que habría resuelto más de ellas. aspeed
\citep{hoos2015aspeed} formula la construcción del \textit{schedule} como un
problema de optimización resuelto con programación de conjuntos de
respuestas.

SUNNY \citep{amadini2014sunny}, la referencia metodológica de este trabajo,
simplifica la idea de CPHydra sin fase de entrenamiento. Para una instancia
nueva busca sus $k$ vecinos más cercanos en una base de conocimiento y
selecciona el menor subconjunto de \textit{solvers} que resuelve el máximo
número de ellos. Divide el tiempo en franjas iguales y asigna a cada
\textit{solver} tantas como vecinos resuelve. Las franjas de los vecinos que
ninguno resuelve pasan a un \textit{solver} de respaldo, y el orden de
ejecución sigue el tiempo medio en el vecindario. Para optimización, SUNNY-CP
\citep{amadini2015sunnycp} reemplaza el conteo de instancias resueltas por una
puntuación que premia probar el óptimo y la calidad de la solución.

%------------------------------------------------------------------------------
\section{Aprendizaje automático para predecir el desempeño de algoritmos}
\label{sec:ml}
%------------------------------------------------------------------------------

Los modelos que predicen el desempeño de un algoritmo pueden ser de
regresión, si estiman una magnitud continua como el tiempo de ejecución, o de
clasificación, si predicen una categoría como ``resuelve antes del límite''.
Predecir el tiempo a partir de características es una técnica consolidada
\citep{hutter2014}. Entre los modelos más usados están los conjuntos de
árboles, tales como el \textbf{bosque aleatorio} \citep{breiman2001}, que
promedia muchos árboles entrenados sobre muestras aleatorias y es uno de los
más precisos para esta tarea \citep{hutter2014}, y el \textit{gradient
boosting} \citep{friedman2001}, que construye los árboles en secuencia, cada
uno ajustado a los errores de los anteriores. Los tiempos suelen estar
\textbf{censurados}, porque cuando un algoritmo agota el límite solo se sabe
que su tiempo real es mayor, por lo cual conviene predecir su logaritmo y
tratar con cuidado esos valores censurados \citep{hutter2014}.

El método de $k$ vecinos más cercanos ($k$-NN) \citep{cover1967} predice el
comportamiento de una instancia a partir de las $k$ más parecidas según una
distancia en el espacio de características. Como su ``modelo'' es la propia
base de instancias, es sensible a la escala de las características, que se
normalizan antes de calcular distancias, y a las características
irrelevantes o redundantes. SUNNY usa $k$-NN para construir el vecindario, y
en sus experimentos valores de $k$ entre 5 y 20 dan resultados similares
\citep{amadini2014sunny}.

\subsection{Métricas para evaluar portafolios}
\label{sec:metricas_portafolios}

Un portafolio se evalúa frente a dos referencias \citep{kerschke2019}. El
\textbf{SBS} (\textit{Single Best Solver}) es el algoritmo con mejor
desempeño promedio, lo que se obtendría usando siempre el mismo, de modo que un
portafolio solo aporta si logra superarlo. El \textbf{VBS} (\textit{Virtual
Best Solver}) es un oráculo que elige en cada instancia el mejor algoritmo, y
marca el límite de lo que puede lograr cualquier portafolio con esos
algoritmos. La diferencia entre ambos mide el margen disponible, y una medida
habitual es la fracción de ese margen que el portafolio logra cerrar
\citep{kerschke2019, lindauer2019}.

Las métricas de desempeño más usadas son el número de instancias resueltas
(en optimización, con el óptimo probado), el \textbf{PAR10}, que promedia el
tiempo contando cada instancia no resuelta como $10 \cdot T$
\citep{lindauer2019}, el tiempo hasta probar el óptimo y el
\textbf{\textit{gap}} frente a la mejor solución conocida
(Sección~\ref{sec:bks_gap}), que mide la calidad cuando el óptimo no se
prueba. SUNNY-CP usa además una puntuación que vale 1 si se prueba el óptimo,
entre 0.25 y 0.75 según la calidad de la mejor solución y 0 si no se
encuentra ninguna \citep{amadini2015sunnycp}.

%------------------------------------------------------------------------------
\section{Estado del arte}
\label{sec:estado_arte}
%------------------------------------------------------------------------------

El estado del arte se construyó mediante una revisión sistemática de
literatura que siguió la guía de \citet{carrera2022}. El protocolo PICOC
(población, intervención, comparación, resultado y contexto) delimitó el
alcance de la búsqueda (Tabla~\ref{tab:picoc}).

\begin{table}[H]
  \centering
  \small
  \singlespacing
  \caption{Protocolo PICOC de la revisión sistemática.}
  \label{tab:picoc}
  \begin{tabularx}{\textwidth}{>{\raggedright\arraybackslash}p{0.16\textwidth}
                               >{\raggedright\arraybackslash}p{0.22\textwidth}
                               >{\raggedright\arraybackslash}X}
    \toprule
    \textbf{Criterio} & \textbf{Término principal} & \textbf{Términos relacionados} \\
    \midrule
    Población    & Job Shop Scheduling Problem & JSSP, \textit{makespan},
      \textit{scheduling} NP-duro. Se restringe a la versión clásica. \\
    Intervención & Portafolio de algoritmos & \textit{Algorithm portfolio},
      \textit{schedule} de ejecución, \textit{Winner Takes All}. \\
    Comparación  & Algoritmo único & Búsqueda tabú, algoritmos genéticos,
      recocido simulado, programación lineal entera mixta. \\
    Resultado    & Selección automática & Problema de selección de
      algoritmos, características de instancias, selección por instancia. \\
    Contexto     & \textit{Machine Learning} & Aprendizaje supervisado,
      bosques aleatorios, máquinas de soporte vectorial, redes neuronales. \\
    \bottomrule
  \end{tabularx}
\end{table}

Del protocolo se derivaron tres preguntas de investigación, que guiaron la
extracción de información de cada artículo:

\begin{description}
  \item[RQ1.] ¿Cuáles son los enfoques de portafolios de algoritmos más
    efectivos aplicados al JSSP, comparados con la estrategia
    \textit{Winner Takes All}?
  \item[RQ2.] ¿Qué características de las instancias del JSSP son más
    relevantes para predecir el desempeño de los algoritmos de resolución?
  \item[RQ3.] ¿Cómo se aplican las técnicas de \textit{Machine Learning} para
    construir \textit{schedules} de ejecución de portafolios en problemas de
    optimización combinatoria NP-duros?
\end{description}

La búsqueda priorizó bases de datos y repositorios de acceso libre, entre
ellos arXiv y dblp, y se limitó a artículos en inglés publicados entre 2001 y 2025 en revistas o conferencias
con revisión por pares y relacionados con al menos dos criterios PICOC. Se
excluyeron los libros sin acceso digital, los trabajos en otros idiomas que no
fueran inglés o español y los centrados exclusivamente en el Flexible Job Shop
o en variantes multiobjetivo. La revisión seleccionó 12 trabajos, organizados
en tres bloques temáticos (Tabla~\ref{tab:slr}). Durante el desarrollo se
incorporaron otras referencias, identificadas a partir de estos trabajos y de
las necesidades de cada fase, sin aplicar de nuevo el protocolo.

\begin{table}[H]
  \centering
  \small
  \singlespacing
  \caption{Trabajos seleccionados en la revisión sistemática.}
  \label{tab:slr}
  \begin{tabularx}{\textwidth}{>{\raggedright\arraybackslash}p{0.30\textwidth}
                               >{\raggedright\arraybackslash}p{0.17\textwidth}
                               >{\raggedright\arraybackslash}X}
    \toprule
    \textbf{Trabajo} & \textbf{Bloque} & \textbf{Aporte al proyecto} \\
    \midrule
    \citet{pinedo2022} & JSSP & Marco formal del JSSP, \textit{makespan} y
      métodos de solución. \\
    \citet{reijnen2023} & JSSP & Plataforma unificada de instancias y métodos
      de referencia para el JSSP y sus variantes. \\
    \citet{xu2025learn} & JSSP / ML & Compara programación genética y
      aprendizaje por refuerzo para aprender heurísticas del JSSP. \\
    \citet{smit2024gnn} & JSSP / ML & Revisión de redes neuronales de grafos
      aplicadas al JSSP. \\
    \citet{gomes2001} & Portafolios & Trabajo fundacional sobre portafolios de
      algoritmos y la variabilidad de su desempeño. \\
    \citet{xu2008satzilla} & Portafolios & Portafolio de selección por
      instancia; arquitectura de referencia. \\
    \citet{streeter2008} & Portafolios & Diseño de portafolios y
      \textit{schedules} de ejecución con reinicio o suspensión. \\
    \citet{amadini2014sunny} & Portafolios & \textit{Schedule} de ejecución
      basado en vecinos más cercanos; referencia metodológica del trabajo. \\
    \citet{omahony2008cphydra} & Portafolios & Portafolio con
      \textit{schedule} de ejecución para programación por restricciones,
      basado en las instancias más parecidas; el primero de su tipo según
      \citet{amadini2015sunnycp}. \\
    \citet{kotthoff2014} & Selección & Clasificación de las técnicas de
      aprendizaje para la selección de algoritmos. \\
    \citet{kerschke2019} & Selección & Revisión de la selección automática de
      algoritmos y de la caracterización de instancias. \\
    \citet{bengio2021} & ML en optimización & Marco metodológico para
      integrar aprendizaje automático y optimización combinatoria. \\
    \bottomrule
  \end{tabularx}
\end{table}
% Nota: el informe de la revisión anotó como "Cunha et al. (2025)", "Hao et al.
% (2024)" y "Xu et al. (2012)" los trabajos de Xu et al. (2025), Smit et al.
% (2024) y Streeter y Smith (2008); se citan con sus autores reales. El informe
% cita Kotthoff (2016), versión en libro del survey; aquí se usa Kotthoff (2014).

\subsection{Síntesis de la revisión}
\label{sec:evolucion_jssp}
\label{sec:ml_co}
% Fuente: Anteproyecto, Cap. 6.2.

Los métodos para el JSSP abarcan desde los primeros algoritmos heurísticos
\citep{giffler1960} y los métodos exactos basados en ramificación y
acotamiento y programación entera hasta metaheurísticas y métodos híbridos
que incorporan aprendizaje \citep{pinedo2022}. En paralelo, los portafolios de
algoritmos se consolidaron en satisfacción de restricciones, planificación y
otros problemas NP-duros, con sistemas como SATzilla \citep{xu2008satzilla}
e Hydra \citep{xu2010hydra}, que caracterizan las instancias y predicen el
algoritmo o la configuración más apropiada. \citet{kotthoff2014} describe
portafolios de algoritmos y de configuraciones, y señala la ejecución paralela
como una línea de investigación en curso. En optimización combinatoria, el aprendizaje automático se ha usado para
predecir el rendimiento de los algoritmos, guiar búsquedas heurísticas y
ajustar parámetros de optimización \citep{bengio2021}, y en el JSSP en
particular para aprender heurísticas y reglas de despacho
\citep{smit2024gnn, xu2025learn}.

%------------------------------------------------------------------------------
\section{Síntesis y brecha identificada}
\label{sec:brecha}
%------------------------------------------------------------------------------

De la literatura revisada se desprende que el desempeño de los métodos para el
JSSP varía de manera pronunciada entre instancias. Los portafolios aprovechan
esa variabilidad, aunque la mayoría elige un único algoritmo por instancia y
queda expuesta a errores de selección que los \textit{schedules} secuenciales
como CPHydra, SUNNY o aspeed mitigan repartiendo el tiempo. Asimismo, el
aprendizaje automático, en particular los vecinos más cercanos y los conjuntos
de árboles, es la herramienta habitual para conectar las características de una
instancia con el desempeño esperado de cada algoritmo.

En los trabajos revisados no se encontró un portafolio por instancia ni un
\textit{schedule} de ejecución al estilo de SATzilla o SUNNY aplicado al JSSP
clásico, ya que los portafolios secuenciales de programación por restricciones
se han evaluado sobre colecciones heterogéneas con características genéricas.
Las revisiones sobre aprendizaje automático en el JSSP se concentran en la
generación de heurísticas y políticas de despacho, con escasa atención a la
selección de algoritmos, y tampoco se encontró un conjunto de datos público con
el desempeño de varios métodos sobre el JSSP, análogo a ASlib
\citep{bischl2016aslib}.

Este trabajo aborda esas brechas aplicando y comparando la selección por
instancia y los \textit{schedules} de ejecución en el JSSP clásico, con
\textit{solvers} de restricciones que comparten un modelo en MiniZinc, y
construye una base de conocimiento con el desempeño de cada configuración en
las familias de instancias de referencia del problema
(Capítulo~\ref{cap:caracterizacion}). Con ella estudia si un
\textit{schedule} aporta frente a la selección de una sola configuración
cuando el portafolio se especializa en un solo problema, y qué tan
informativas son las características genéricas y las propias del JSSP.

%==============================================================================
\chapter{Definición formal del problema y metodología}
\label{cap:problema}
%==============================================================================

A partir de los fundamentos anteriores, el diseño del portafolio requiere fijar
la formulación matemática del JSSP en MiniZinc y el esquema formal para
seleccionar y secuenciar las configuraciones, junto con la metodología seguida.

%------------------------------------------------------------------------------
\section{Notación y definición del JSSP}
\label{sec:notacion}
%------------------------------------------------------------------------------

Una instancia del JSSP clásico se define por un conjunto de $n$ trabajos
$J = \{1, \dots, n\}$ y un conjunto de $m$ máquinas $M = \{0, \dots, m-1\}$.
Cada trabajo $j$ es una secuencia ordenada de $m$ operaciones
$o_{j1}, \dots, o_{jm}$, y cada operación $o_{jk}$ debe procesarse en la
máquina $\mu_{jk} \in M$ durante $p_{jk} \in \mathbb{Z}^{+}$ unidades de
tiempo. En la versión clásica cada trabajo pasa exactamente una vez por cada
máquina, de modo que hay $n \cdot m$ operaciones, condición que cumplen todas
las instancias de este trabajo. Las variables de decisión son el tiempo de
inicio $s_{jk} \in \mathbb{Z}_{\ge 0}$ de cada operación y el
\textit{makespan} $C_{\max}$, el instante en que termina la última operación.

Una operación no puede empezar hasta que termine la anterior de su trabajo
(precedencia), cada máquina procesa una sola operación a la vez y sin
interrupciones (capacidad), y el \textit{makespan} no puede ser menor que el
final de ningún trabajo \citep{pinedo2022}:
\begin{align}
  & s_{jk} + p_{jk} \le s_{j,k+1},
    && \forall j \in J,\ k = 1, \dots, m-1, \label{eq:precedencia} \\
  & s_{jk} + p_{jk} \le s_{j'k'} \;\lor\; s_{j'k'} + p_{j'k'} \le s_{jk},
    && \forall\, o_{jk} \ne o_{j'k'} : \mu_{jk} = \mu_{j'k'}, \label{eq:disyuntiva} \\
  & C_{\max} \ge s_{jm} + p_{jm}, && \forall j \in J. \label{eq:makespan}
\end{align}
El objetivo es
\begin{equation}
  \min\; C_{\max}.
  \label{eq:objetivo}
\end{equation}
Las restricciones~\eqref{eq:disyuntiva} se llaman disyuntivas porque expresan
una elección entre dos órdenes posibles, y es esa elección combinatoria la que
hace difícil el problema.

El valor óptimo $C^{*}_{\max}$ admite una cota inferior simple, porque ningún
programa puede terminar antes de que el trabajo más largo complete sus
operaciones ni antes de que la máquina más cargada procese todas las suyas,
la misma cota que \citet{taillard1993} reporta para sus instancias:
\begin{equation}
  C^{*}_{\max} \ge \max\Bigl( \max_{j \in J} \sum_{k=1}^{m} p_{jk},\;
  \max_{r \in M} \sum_{(j,k) : \mu_{jk} = r} p_{jk} \Bigr).
  \label{eq:cota_inferior}
\end{equation}
Por otro lado, ejecutar todas las operaciones una tras otra siempre es
factible, así que $C^{*}_{\max} \le H = \sum_{j,k} p_{jk}$, el horizonte
trivial.

%------------------------------------------------------------------------------
\section{Modelos en MiniZinc}
\label{sec:modelo_minizinc}
%------------------------------------------------------------------------------

El trabajo usa tres modelos de MiniZinc que traducen la misma formulación y,
por tanto, tienen las mismas soluciones óptimas. Se diferencian en la
información adicional que reciben los \textit{solvers}, dando lugar al modelo
original (\texttt{jssp.mzn}), al modelo con cotas (\texttt{jssp\_cotas.mzn}),
que usa el portafolio, y al modelo con cotas y estrategia de búsqueda
(\texttt{jssp\_\allowbreak cotas\_\allowbreak busqueda.mzn}), con el que Chuffed
forma la tercera configuración del portafolio.

\subsection{Modelo original}
\label{sec:modelo_original}

En el modelo original los tiempos de inicio y el \textit{makespan} tienen
dominio $0..H$ y cada \textit{solver} decide cómo buscar. Las instancias se
entregan como archivos \texttt{.dzn} con cuatro parámetros que especifican el
número de trabajos y de máquinas junto con las matrices \texttt{machines}
($\mu_{jk}$) y \texttt{durations} ($p_{jk}$), en las que la fila $j$ describe
las operaciones del trabajo $j$ en su orden de ejecución. La precedencia se
expresa directamente con la ecuación~\eqref{eq:precedencia}. Para la
capacidad, el modelo reúne, en cada máquina, los tiempos de inicio y las
duraciones de las $n$ operaciones que la usan e impone la restricción global
\texttt{disjunctive}, equivalente al conjunto de
disyunciones~\eqref{eq:disyuntiva}, que permite a cada \textit{solver} usar un
propagador especializado. El código completo se presenta en la
Sección~\ref{sec:modelo_cotas}, ya que el modelo con cotas solo modifica
algunas líneas.

El horizonte trivial constituye una cota muy holgada, puesto que en las 242
instancias $H$ es entre 3.6 y 18.8 veces la mejor solución conocida (11.4
veces en promedio). Con dominios tan amplios, un \textit{solver} puede aceptar
como primera solución un programa casi secuencial y avanzar muy lentamente
desde ahí. Este comportamiento, que se observó en Chuffed
(Sección~\ref{sec:ejecucion_solvers}), motivó la incorporación de cotas.
La idea de acotar el modelo surgió de ese diagnóstico y se llevó a cabo con
técnicas conocidas en \textit{scheduling}: la cota inferior de la
ecuación~\eqref{eq:cota_inferior} y el ajuste de las \textit{cabezas} y
\textit{colas} de cada operación respecto a una cota superior del
\textit{makespan} \citep{carlier1994}.

\subsection{Cotas del \textit{makespan}}
\label{sec:cotas}

Antes de resolver una instancia se calculan dos cotas de su óptimo. La cota
inferior $lb$ es la de la ecuación~\eqref{eq:cota_inferior}. La cota superior
$ub$ es el \textit{makespan} de un programa factible construido con el
algoritmo de \citet{giffler1960}, que genera un programa activo, es decir, uno
en el que ninguna operación puede adelantarse sin retrasar otra. El algoritmo
programa una operación por paso:

\begin{enumerate}
  \item entre las operaciones disponibles (la siguiente de cada trabajo), toma
    la que puede terminar más temprano, $o^{*}$, y su máquina $r^{*}$;
  \item forma el conjunto de conflicto con las operaciones disponibles que
    usan $r^{*}$ y pueden empezar antes de que termine $o^{*}$;
  \item elige una de ellas con una regla de prioridad y la programa en su
    inicio más temprano, hasta programar las $n \cdot m$ operaciones.
\end{enumerate}

\noindent Se usaron seis reglas de prioridad clásicas \citep{panwalkar1977}
(Tabla~\ref{tab:reglas}). Para cada instancia se construye un programa con
cada regla, se verifica que cumpla las restricciones de precedencia y de
capacidad y se toma como $ub$ el menor \textit{makespan}. El cálculo de ambas
cotas para las 242 instancias toma pocos segundos
(\texttt{src/02\_generar\_cotas.py}).

\begin{table}[H]
  \centering
  \small
  \singlespacing
  \caption{Reglas de prioridad usadas en el algoritmo de Giffler y Thompson.}
  \label{tab:reglas}
  \begin{tabularx}{\textwidth}{l >{\raggedright\arraybackslash}X r}
    \toprule
    \textbf{Regla} & \textbf{Elige la operación\dots} & \textbf{Veces mejor} \\
    \midrule
    SPT  & de menor duración & 0 \\
    LPT  & de mayor duración & 0 \\
    MWKR & del trabajo con más tiempo de proceso pendiente & 91 \\
    LWKR & del trabajo con menos tiempo de proceso pendiente & 3 \\
    MOR  & del trabajo con más operaciones pendientes & 58 \\
    FIFO & que quedó disponible primero & 90 \\
    \bottomrule
  \end{tabularx}
  \par\vspace{2pt}{\footnotesize Veces mejor: número de instancias en que la regla
  dio la menor $ub$ (en caso de empate se cuenta la primera regla de la lista).}
\end{table}

En promedio, la cota inferior equivale al 88.3\,\% de la mejor solución
conocida y la cota superior la supera en un 26.3\,\% (detalle por familia en
la Tabla~\ref{tab:cotas}). En seis instancias (\texttt{la06}, \texttt{la10},
\texttt{la12}, \texttt{la14}, \texttt{swv16} y \texttt{swv18}) ambas cotas
coinciden, de modo que el programa heurístico es óptimo y su optimalidad queda
probada sin ejecutar ningún \textit{solver}.

\subsection{Modelo con cotas}
\label{sec:modelo_cotas}

Un aporte realizado en este trabajo es integrar, mediante técnicas habituales
de reducción de dominios, las cotas del preprocesamiento de Giffler y Thompson
en un modelo con cotas en MiniZinc que recibe $lb$ y $ub$ como parámetros: los tiempos de inicio pasan a $0..ub$, el
\textit{makespan} queda en $lb..ub$ y cada operación no puede empezar antes de
las anteriores de su trabajo ni terminar después de $ub$ restadas las
siguientes, lo que corresponde a acotar sus cabezas y colas \citep{carlier1994},
\begin{equation}
  \sum_{l=1}^{k-1} p_{jl} \;\le\; s_{jk} \;\le\; ub - \sum_{l=k}^{m} p_{jl},
  \qquad \forall j \in J,\ k = 1, \dots, m.
  \label{eq:cabeza_cola}
\end{equation}
El código del modelo, sin la sección de salida, se presenta a continuación.
Las líneas marcadas con \texttt{(*)} son las que lo distinguen del modelo
original, en el que \texttt{s} y \texttt{makespan} tienen dominio
\texttt{0..horizon}, con \texttt{horizon = sum(durations)}:

\noindent\begin{minipage}{\linewidth}\singlespacing\small
\begin{verbatim}
include "globals.mzn";
int: n_jobs;  int: n_machines;
array[1..n_jobs, 1..n_machines] of int: machines;
array[1..n_jobs, 1..n_machines] of int: durations;
int: lb;  int: ub;                                              % (*)
array[1..n_jobs, 1..n_machines] of var 0..ub: s;                % (*)
var lb..ub: makespan;                                           % (*)

constraint forall(j in 1..n_jobs, o in 1..n_machines-1) (
    s[j, o] + durations[j, o] <= s[j, o+1]);
constraint forall(m in 0..n_machines-1) (
    let { array[1..n_jobs] of var int: ms =
            [s[j, o] | j in 1..n_jobs, o in 1..n_machines
             where machines[j, o] == m];
          array[1..n_jobs] of int: md =
            [durations[j, o] | j in 1..n_jobs, o in 1..n_machines
             where machines[j, o] == m];
    } in disjunctive(ms, md));
constraint forall(j in 1..n_jobs) (
    s[j, n_machines] + durations[j, n_machines] <= makespan);
constraint forall(j in 1..n_jobs, o in 1..n_machines) (        % (*)
    s[j, o] >= sum(k in 1..o-1)(durations[j, k]) /\
    s[j, o] + sum(k in o..n_machines)(durations[j, k]) <= ub);
solve minimize makespan;
\end{verbatim}
\end{minipage}
\vspace{\baselineskip}

Como $lb \le C^{*}_{\max} \le ub$, estas restricciones conservan todas las
soluciones óptimas y descartan únicamente los programas peores que el
heurístico. Con ellas los dominios se reducen de forma pronunciada y la
propagación es más efectiva desde el inicio de la búsqueda. A cambio, la
primera solución que encuentre un \textit{solver} ya debe ser al menos tan
buena como la heurística, lo que en instancias grandes puede ser difícil para
un \textit{solver} sin aprendizaje. Por eso el portafolio conserva siempre el
programa heurístico, de manera que si un \textit{solver} no encuentra
solución, el resultado es la $ub$.

\subsection{Modelo con estrategia de búsqueda}
\label{sec:modelo_busqueda}

Otro aporte de este trabajo es obtener del mismo \textit{solver} una
configuración adicional del portafolio (Chuffed V2) mediante una anotación de
búsqueda estándar de MiniZinc, la
cual fija en cada paso la operación con menor inicio posible en su valor
mínimo. Para ello reemplaza la última línea por

{\singlespacing\small
\begin{verbatim}
solve :: int_search([s[j, o] | j in 1..n_jobs, o in 1..n_machines]
                    ++ [makespan], smallest, indomain_min) minimize makespan;
\end{verbatim}
}

\noindent y se ejecuta sin la opción \texttt{-f}, con lo cual el
\textit{solver} aplica esta heurística de ramificación. La búsqueda libre de
Chuffed es muy eficaz para probar la optimalidad en instancias pequeñas y
medianas, pero en las grandes puede tardar en encontrar una primera solución
mejor que la heurística. La estrategia explícita favorece encontrar buenas
soluciones pronto en instancias grandes, a cambio de probar menos óptimos, por
lo que ambas variantes son candidatas naturales a complementarse dentro del
portafolio (Sección~\ref{sec:composicion}).

%------------------------------------------------------------------------------
\section{Formulación del problema de selección y del \textit{schedule}}
\label{sec:formulacion_schedule}
%------------------------------------------------------------------------------

\subsection{Portafolio, instancias y presupuesto de tiempo}
Sea $\Pi = \{a_1, \dots, a_q\}$ el portafolio, $\mathcal{I}$ el conjunto de
instancias y $T$ el presupuesto de tiempo por instancia (300 segundos). Cada
elemento del portafolio es una \textit{configuración}, entendida como un
\textit{solver} junto con el modelo que resuelve. En este trabajo, $\Pi$ tiene
tres configuraciones conformadas por Chuffed y Gecode con el modelo con cotas,
y por Chuffed con el modelo con estrategia de búsqueda (en adelante, Chuffed
V2). Para cada configuración $a \in \Pi$ e instancia $i \in \mathcal{I}$, la
base de conocimiento registra el tiempo $t(a, i)$ que tarda $a$ en probar el
óptimo de $i$ (con $t(a,i) = \infty$ si no lo prueba dentro de $T$) y el mejor
\textit{makespan} $C(a, i, \tau)$ encontrado por $a$ en $i$ tras $\tau$
segundos de ejecución. La función $C(a, i, \tau)$ se obtiene de la
\textit{trayectoria} de la ejecución, es decir, la lista de soluciones
mejoradas que reporta el \textit{solver}, cada una con el instante en que se
encontró. Si $a$ todavía no ha encontrado ninguna solución en el instante
$\tau$, se toma $C(a, i, \tau) = ub(i)$, porque el portafolio dispone siempre
del programa heurístico. Se dice que $a$ resuelve $i$ en un tiempo $\tau$ si
$t(a, i) \le \tau$. Cada instancia se representa además por su vector de
características $f(i) \in \mathbb{R}^{d}$ (Sección~\ref{sec:features}).

Las instancias con $lb(i) = ub(i)$ quedan resueltas por las cotas antes de
ejecutar cualquier configuración. Las comparaciones entre métodos se hacen
sobre las instancias restantes.

\subsection{Selección de un único algoritmo}
En la estrategia \textit{Winner Takes All}, un selector
$S: \mathbb{R}^{d} \rightarrow \Pi$ elige una configuración a partir de las
características de la instancia, y esa configuración se ejecuta durante todo
el presupuesto. El resultado sobre $i$ es el de $S(f(i))$ con tiempo $T$, por
lo cual la instancia se resuelve si $t(S(f(i)), i) \le T$, y la calidad
obtenida es $C(S(f(i)), i, T)$.

Con esta notación se definen las dos referencias de la
Sección~\ref{sec:metricas_portafolios}. El \textbf{SBS} es el selector
constante que elige siempre la configuración con mejor desempeño medio en
$\mathcal{I}$. El \textbf{VBS} es el selector ideal que, para cada instancia,
elige la configuración con mejor desempeño en ella, definida bajo el criterio
de calidad de este trabajo como $\arg\min_{a \in \Pi} C(a, i, T)$.

\subsection{\textit{Schedule} de ejecución}
Un \textit{schedule} de ejecución para la instancia $i$ es una secuencia
\begin{equation}
  \sigma(i) = \bigl[ (a_{(1)}, \tau_1), (a_{(2)}, \tau_2), \dots,
  (a_{(r)}, \tau_r) \bigr],
  \qquad \sum_{l=1}^{r} \tau_l \le T,
  \label{eq:schedule}
\end{equation}
donde $a_{(l)} \in \Pi$ son configuraciones, que pueden repetirse en
ejecuciones independientes, y $\tau_l > 0$ el
tiempo asignado a cada una. Las configuraciones se ejecutan en el orden
indicado, cada una durante su tiempo y de forma independiente, sin transmitir
la mejor solución de una a la siguiente. El resultado del \textit{schedule} es
\begin{equation}
  C(\sigma, i) = \min_{l = 1, \dots, r} C\bigl(a_{(l)}, i, \tau_l\bigr).
  \label{eq:resultado_schedule}
\end{equation}
La instancia se resuelve si alguna configuración la resuelve en su
intervalo, es decir, si existe $l$ tal que $t(a_{(l)}, i) \le \tau_l$, y el
tiempo empleado es la suma de los intervalos anteriores más
$t(a_{(l)}, i)$. Gracias a las trayectorias, la
ecuación~\eqref{eq:resultado_schedule} permite evaluar cualquier
\textit{schedule} sin volver a ejecutar los \textit{solvers}.

Construir un \textit{schedule} implica decidir qué configuraciones incluir
(el subconjunto $\{a_{(1)}, \allowbreak \dots, \allowbreak a_{(r)}\} \subseteq \Pi$), cuánto tiempo
darle a cada una (los valores $\tau_l$) y en qué orden ejecutarlas, que son
las tres decisiones de la pregunta de investigación. La estrategia
\textit{Winner Takes All} es el caso particular $r = 1$ y $\tau_1 = T$. En
este trabajo se comparan varias formas de tomar estas decisiones, que abarcan
un \textit{schedule} fijo para todas las instancias, el enfoque de SUNNY
\citep{amadini2014sunny}, basado en las instancias más parecidas, y un
\textit{schedule} construido con modelos de aprendizaje automático. Su
implementación se describe en el Capítulo~\ref{cap:portafolio}.

%------------------------------------------------------------------------------
\section{Metodología de trabajo}
\label{sec:metodologia}
%------------------------------------------------------------------------------
% Fuente: Anteproyecto, Cap. 7.1.

El desarrollo del proyecto se estructuró en cinco fases, cuyo desglose de
actividades se recoge en la Tabla~\ref{tab:fases}.

\begin{table}[H]
  \centering
  \small
  \singlespacing
  \caption{Fases y actividades del trabajo.}
  \label{tab:fases}
  \begin{tabularx}{\textwidth}{>{\raggedright\arraybackslash}p{0.33\textwidth}
                               >{\raggedright\arraybackslash}X}
    \toprule
    \textbf{Fase} & \textbf{Actividades} \\
    \midrule
    1. Diseñar e implementar un portafolio de algoritmos de optimización para
    el JSSP
      & A1. Identificación y selección de los algoritmos a incluir en el
          portafolio.\newline
        A2. Diseño de las adaptaciones necesarias de los algoritmos para el
          JSSP.\newline
        A3. Implementación de los algoritmos seleccionados.\newline
        A4. Validación del correcto funcionamiento de los algoritmos
          implementados. \\
    \addlinespace
    2. Desarrollar un sistema de caracterización de instancias del problema
      & A5. Identificación de las características relevantes para describir
          instancias del JSSP.\newline
        A6. Diseño del sistema de extracción automática de
          características.\newline
        A7. Implementación del sistema de extracción automática de
          características. \\
    \addlinespace
    3. Construir el conjunto de datos para el entrenamiento del sistema de
    selección
      & A8. Selección de instancias representativas para conformar el conjunto
          de datos.\newline
        A9. Ejecución de los algoritmos sobre las instancias
          seleccionadas.\newline
        A10. Preprocesamiento y preparación del conjunto de datos final. \\
    \addlinespace
    4. Desarrollar el modelo de selección automática de algoritmos
      & A11. Diseño del \textit{pipeline} de \textit{machine learning} para la
          selección de algoritmos.\newline
        A12. Entrenamiento y optimización de los modelos de
          predicción.\newline
        A13. Selección del modelo final con base en métricas de desempeño. \\
    \addlinespace
    5. Implementar el sistema integrado de portafolios de algoritmos
      & A14. Diseño de la arquitectura del sistema integrado.\newline
        A15. Implementación del sistema completo.\newline
        A16. Evaluación experimental y análisis de resultados. \\
    \bottomrule
  \end{tabularx}
\end{table}

Durante el desarrollo, y con la aprobación de la dirección del trabajo, se
hicieron tres ajustes que conservan intactos los objetivos trazados:

\begin{itemize}
  \item En las actividades A2 y A9, el análisis de la primera campaña mostró
    que el horizonte trivial limitaba a Chuffed en las instancias grandes. Se
    incorporaron las cotas de la Sección~\ref{sec:cotas} y se repitió la
    campaña con el modelo acotado, donde la heurística de Giffler y Thompson
    calcula la cota superior y entrega una solución de respaldo.
  \item En las actividades A1 y A2, el portafolio se compuso de tres
    configuraciones basadas en Chuffed y Gecode con el modelo con cotas y en
    Chuffed con estrategia de búsqueda, siguiendo la idea de portafolios de
    configuraciones como Hydra \citep{xu2010hydra}. CP-SAT se empleó como
    referencia externa, dado que domina a los otros dos en el modelo original y
    acapararía la selección (Sección~\ref{sec:composicion}).
  \item En las actividades A11 a A13, la construcción del \textit{schedule} se
    apoyó en un esquema híbrido compuesto por un clasificador de optimalidad
    probada, un regresor del tiempo condicionado a probarla y un clasificador
    de la mejor configuración (Sección~\ref{sec:modelos_prediccion}).
\end{itemize}

%==============================================================================
\chapter{Caracterización de instancias y conjunto de datos}
\label{cap:caracterizacion}
%==============================================================================
% Objetivo específico 2. Actividades A5 a A10.

Para llevar a la práctica los modelos de MiniZinc y la formulación del
\textit{schedule} definidos en el capítulo anterior, es necesario
reunir un conjunto de instancias de prueba, extraer sus descriptores y medir
en ellas el comportamiento real de cada configuración. Ese proceso abarca la
conversión de las instancias (Sección~\ref{sec:instancias}), el cálculo de sus
cotas (Sección~\ref{sec:cotas_instancias}), la extracción de características
(Sección~\ref{sec:features}) y las campañas de ejecución de los
\textit{solvers} (Sección~\ref{sec:ejecucion_solvers}), cuyos registros se
consolidan en la base de conocimiento (Sección~\ref{sec:preprocesamiento}) que
alimenta las estrategias de selección del Capítulo~\ref{cap:portafolio}.

%------------------------------------------------------------------------------
\section{Conjunto de instancias}
\label{sec:instancias}
%------------------------------------------------------------------------------

\subsection{Familias de instancias}

Se trabajó con 242 instancias clásicas del JSSP, distribuidas en ocho familias
ampliamente utilizadas como \textit{benchmark} en la literatura. Varias de ellas
se distribuyen a través de la OR-Library \citep{beasley1990}, y todas están
reunidas, junto con sus cotas y mejores soluciones conocidas, en el repositorio
de \citet{weise2019jssp}. En la Tabla~\ref{tab:familias} se detallan el origen,
el número de instancias y las dimensiones de cada familia.

\begin{table}[H]
  \centering
  \small
  \singlespacing
  \caption{Familias de instancias utilizadas. El tamaño se expresa como
  $n \times m$ (trabajos $\times$ máquinas).}
  \label{tab:familias}
  \begin{tabularx}{\textwidth}{>{\ttfamily}l >{\raggedright\arraybackslash}X
                               c >{\raggedright\arraybackslash}p{0.33\textwidth}}
    \toprule
    \normalfont\textbf{Familia} & \textbf{Origen} & \textbf{Instancias} & \textbf{Tamaños} \\
    \midrule
    ft  & \citet{fisher1963}         & 3  & $6\times6$, $10\times10$, $20\times5$ \\
    la  & \citet{lawrence1984}       & 40 & de $10\times5$ a $30\times10$ \\
    abz & \citet{adams1988}          & 5  & $10\times10$, $20\times15$ \\
    orb & \citet{applegate1991}      & 10 & $10\times10$ \\
    swv & \citet{storer1992}         & 20 & $20\times10$, $20\times15$, $50\times10$ \\
    yn  & \citet{yamada1992}         & 4  & $20\times20$ \\
    ta  & \citet{taillard1993}       & 80 & de $15\times15$ a $100\times20$ \\
    dmu & \citet{demirkol1998}       & 80 & de $20\times15$ a $50\times20$ \\
    \midrule
    \normalfont\textbf{Total} & & 242 & de $6\times6$ a $100\times20$ \\
    \bottomrule
  \end{tabularx}
  \par\vspace{2pt}{\footnotesize El origen de las familias \texttt{ft}, \texttt{la},
  \texttt{abz}, \texttt{orb}, \texttt{swv} e \texttt{yn} es el que indica la
  OR-Library \citep{orlibjobshop}.}
\end{table}

El conjunto cubre un rango amplio de dificultad. Las instancias pequeñas de las
familias \texttt{ft}, \texttt{la} y \texttt{orb} se resuelven a optimalidad en
segundos con métodos exactos modernos, mientras que muchas instancias de
\texttt{ta}, \texttt{dmu}, \texttt{swv} e \texttt{yn} siguen abiertas, ya que
su óptimo no se conoce y solo se dispone de una cota inferior y de la mejor
solución encontrada hasta ahora. De las 242 instancias, 156 tienen óptimo
conocido (su cota inferior publicada es igual a su mejor solución conocida).
Esta diversidad es deseable para un portafolio, porque el desempeño relativo
de los \textit{solvers} suele cambiar con el tamaño y la estructura de la
instancia.

\subsection{Conversión al formato de MiniZinc}

Las instancias se encuentran en el formato de texto de la OR-Library, compuesto
por una primera línea con $n$ y $m$ seguida de una línea por trabajo con $m$
pares \textit{(máquina, duración)} en orden, con las máquinas numeradas desde
cero. El programa \texttt{src/01\_convertir\_instancias.py} las transforma en
archivos \texttt{.dzn} con los parámetros del modelo
(Sección~\ref{sec:modelo_original}), sin modificar los índices, y genera una
segunda versión con los parámetros \texttt{lb} y \texttt{ub}
(Sección~\ref{sec:cotas_instancias}). Se verificó que las matrices de las 242
instancias tuvieran las dimensiones declaradas y que los índices de máquina
estuvieran en el rango $0,\dots,m-1$, y, como prueba de extremo a extremo,
todas las configuraciones resolvieron \texttt{ft06} y obtuvieron su óptimo
conocido, 55.

%------------------------------------------------------------------------------
\section{Cotas de las instancias}
\label{sec:cotas_instancias}
%------------------------------------------------------------------------------

Para cada instancia se calcularon la cota inferior $lb$ y la cota superior
$ub$ de la Sección~\ref{sec:cotas} con el programa
\texttt{src/02\_generar\_cotas.py}. La calidad alcanzada en cada familia frente
a la mejor solución conocida (BKS, Sección~\ref{sec:bks_gap}) se reporta en la
Tabla~\ref{tab:cotas}.

\begin{table}[H]
  \centering
  \small
  \singlespacing
  \caption{Calidad de las cotas por familia.}
  \label{tab:cotas}
  \begin{tabular}{lrrrrrl}
    \toprule
    Familia & Inst. & lb/BKS (\%) & Gap de ub (\%) & $H$/BKS & lb $=$ ub & Regla más frecuente \\
    \midrule
    ft & 3 & 84.0 & 23.1 & 4.5 & 0 & MWKR \\
    la & 40 & 93.5 & 14.7 & 6.6 & 4 & MWKR \\
    abz & 5 & 80.4 & 23.1 & 9.2 & 0 & FIFO \\
    orb & 10 & 70.3 & 26.5 & 5.6 & 0 & FIFO \\
    swv & 20 & 88.1 & 32.4 & 8.3 & 2 & FIFO \\
    yn & 4 & 77.0 & 25.4 & 13.0 & 0 & FIFO \\
    ta & 80 & 90.8 & 22.8 & 13.9 & 0 & FIFO \\
    dmu & 80 & 86.7 & 34.4 & 13.1 & 0 & MWKR \\
    \midrule
    Total & 242 & 88.3 & 26.3 & 11.4 & 6 & MWKR \\
    \bottomrule
  \end{tabular}
  \par\vspace{2pt}{\footnotesize lb: cota inferior; ub: mejor programa de Giffler-Thompson; $H$: horizonte del modelo original (suma de duraciones).}
\end{table}


Las cotas reducen el espacio de búsqueda en un orden de magnitud, puesto que el
horizonte trivial del modelo original es, en promedio, 11.4 veces la mejor
solución conocida, mientras que la cota superior la supera solo en un
26.3\,\%. El efecto es mayor en las familias de instancias grandes
(\texttt{ta}, \texttt{dmu}, \texttt{yn}), donde el horizonte trivial es más de
13 veces la BKS. La heurística ya es, además, una solución razonable, de modo
que el portafolio parte en el peor caso de un programa factible a una
distancia moderada de la mejor solución conocida. Las seis instancias en que
$lb = ub$ no requieren ejecutar ningún \textit{solver} y se excluyen de la
comparación entre estrategias, que se hace sobre las 236 restantes. El mejor
resultado se reparte entre varias reglas de prioridad, ya que MWKR, FIFO y MOR
obtienen la menor cota en distintas familias, lo que justifica evaluar las seis
reglas por instancia.

%------------------------------------------------------------------------------
\section{Extracción de características}
\label{sec:features}
%------------------------------------------------------------------------------

Cada instancia se describe mediante las características genéricas que calcula
\texttt{mzn2feat} a partir del modelo MiniZinc y un grupo de características
propias del JSSP calculadas directamente sobre los datos de la instancia.

\subsection{Características genéricas con \texttt{mzn2feat}}

Se usó \texttt{mzn2feat-1.0} \citep{amadini2014mzn2feat, amadini2015sunnycp},
el extractor de
características del portafolio SUNNY-CP, que traduce el modelo y los datos a
FlatZinc y calcula 95 características numéricas sobre el resultado. Su uso en
portafolios de programación por restricciones está documentado, lo que hace
comparables los resultados con los de SUNNY, y sus características se calculan
exclusivamente a partir de la estructura del modelo que recibe el
\textit{solver}. Como todas las instancias usan el mismo modelo, las
diferencias entre sus vectores se deben solo a los datos. El programa
\texttt{src/03\_caracteristicas\_mzn2feat.py} ejecutó la extracción sobre el
modelo original dentro del contenedor Docker oficial de SUNNY-CP
(\texttt{jacopomauro/sunny-cp}), que incluye la herramienta compilada y
MiniZinc~1.6, ya que \texttt{mzn2feat} se distribuye para versiones antiguas
de MiniZinc. Se usó el modelo original porque el propósito es describir la
instancia, y las cotas se derivan de esos mismos datos. Las características se
agrupan en seis categorías, identificadas por el prefijo de su nombre
(Tabla~\ref{tab:categorias_features}), y dentro de cada una se calculan
estadísticos como la media (\texttt{avg}), el coeficiente de variación
(\texttt{cv}) o la entropía (\texttt{ent}) del grado o del tamaño de dominio
de variables y restricciones (Anexo~\ref{anx:caracteristicas}).

\begin{table}[H]
  \centering
  \small
  \singlespacing
  \caption{Categorías de características de \texttt{mzn2feat} y número de
  características activas tras el preprocesamiento
  (Sección~\ref{sec:limpieza}).}
  \label{tab:categorias_features}
  \begin{tabularx}{\textwidth}{>{\ttfamily}l >{\raggedright\arraybackslash}X c c}
    \toprule
    \normalfont\textbf{Prefijo} & \textbf{Describe} & \textbf{Extraídas} & \textbf{Activas} \\
    \midrule
    v\_  & Variables: número, grado (en cuántas restricciones aparece cada una) y tamaño de dominio & 27 & 26 \\
    d\_  & Tipos de dominio: proporción de variables y restricciones enteras, booleanas, de conjunto y de punto flotante & 18 & 8 \\
    c\_  & Restricciones: número, grado y tamaño de los dominios de sus variables & 27 & 21 \\
    gc\_ & Restricciones globales: cuántas hay y de qué tipo & 4 & 3 \\
    o\_  & Función objetivo: dominio y grado de la variable a optimizar & 8 & 7 \\
    s\_  & Anotaciones de búsqueda incluidas en el modelo & 11 & 0 \\
    \midrule
    \normalfont\textbf{Total} & & 95 & 65 \\
    \bottomrule
  \end{tabularx}
\end{table}

Las 11 características de búsqueda (\texttt{s\_}) son constantes, porque el
modelo original no incluye anotaciones de búsqueda, y lo mismo ocurre con las
que miden tipos de variables o restricciones que el modelo no usa. Al
emplearse un único modelo, muchas de las restantes dependen sobre todo del
tamaño de la instancia. Por ejemplo, el dominio de las variables de inicio es
$0..H$, de modo que los estadísticos de tamaño de dominio crecen con el número
de operaciones y con sus duraciones. Pocas describen la estructura propia del
JSSP, como la distribución de la carga entre máquinas o el orden de las rutas.

\subsection{Características propias del JSSP}
\label{sec:features_propias}

Otro aporte de este trabajo es el diseño de un conjunto de 20 características
propias del JSSP para complementar a \texttt{mzn2feat}, las cuales se calculan
directamente sobre las matrices de máquinas y duraciones de cada instancia con
el programa \texttt{src/\allowbreak 04\_\allowbreak caracteristicas\_\allowbreak propias.py}.
Siguiendo la línea de los trabajos que caracterizan instancias del JSSP a
partir de su estructura \citep{mirshekarian2016}, estos descriptores capturan
el tamaño, la distribución de las duraciones, el equilibrio de la carga entre
máquinas y trabajos, y la forma de las rutas
(Tabla~\ref{tab:features_propias}). Su cálculo toma menos de un segundo para
las 242 instancias y no requiere MiniZinc.

\begin{table}[H]
  \centering
  \small
  \singlespacing
  \caption{Características propias del JSSP. $P$: duraciones; $L_k$: carga
  (suma de duraciones) de la máquina $k$; $W_j$: duración total del trabajo
  $j$; $LB$: cota inferior.}
  \label{tab:features_propias}
  \begin{tabularx}{\textwidth}{l >{\ttfamily\raggedright\arraybackslash}p{4.6cm} >{\raggedright\arraybackslash}X}
    \toprule
    \textbf{Grupo} & \normalfont\textbf{Característica} & \textbf{Definición} \\
    \midrule
    Tamaño & jf\_n, jf\_m & Número de trabajos y de máquinas \\
     & jf\_ratio\_nm, jf\_ops & $n/m$ y número de operaciones $n \cdot m$ \\
    \midrule
    Duraciones & jf\_p\_mean & Duración media de las operaciones \\
     & jf\_p\_cv & Coeficiente de variación de las duraciones \\
     & jf\_p\_range & $(\max P - \min P) / \overline{P}$ \\
    \midrule
    Carga & jf\_load\_cv & Coeficiente de variación de $L_k$ \\
     & jf\_bottleneck & $\max L_k / \overline{L}$: intensidad del cuello de botella \\
     & jf\_job\_cv & Coeficiente de variación de $W_j$ \\
     & jf\_jobmax\_rel & $\max W_j / \overline{W}$ \\
     & jf\_machine\_vs\_job & $\max L_k / \max W_j$: qué componente domina la cota inferior \\
     & jf\_lb\_over\_meanload & $LB / \overline{L}$ \\
     & jf\_lb\_over\_horizon & $LB / H$: holgura del horizonte trivial \\
     & jf\_mach\_pmean\_cv & Coeficiente de variación de la duración media por máquina \\
    \midrule
    Rutas & jf\_route\_pos\_entropy & Entropía media de las máquinas que ocupan cada posición de las rutas \\
     & jf\_route\_kendall & Concordancia media (tipo $\tau$ de Kendall) entre los órdenes de máquinas de dos trabajos; vale 1 en un \textit{flow shop} \\
     & jf\_first\_mach\_ent, jf\_last\_mach\_ent & Entropía de la primera y de la última máquina de las rutas \\
     & jf\_bneck\_pos & Posición media, en las rutas, de las operaciones de la máquina más cargada \\
    \bottomrule
  \end{tabularx}
\end{table}

En la validación del Capítulo~\ref{cap:resultados} se comparan tres conjuntos
conformados por las 65 características activas de \texttt{mzn2feat}, las 20
propias y la unión de ambos (85 características). La última característica de
la tabla resultó redundante, ya que cuando cada máquina procesa exactamente $n$
operaciones, como ocurre en todas las instancias clásicas, coincide con
\texttt{jf\_load\_cv}. Se conservó para mantener el conjunto tal como se
definió antes de la evaluación, con el único efecto de dar un peso doble a esa
magnitud en la distancia de los métodos de vecinos más cercanos.

%------------------------------------------------------------------------------
\section{Ejecución de los solvers}
\label{sec:ejecucion_solvers}
%------------------------------------------------------------------------------

Se realizaron dos campañas de ejecución de 300 segundos por instancia y
configuración. La primera, con el modelo original, sirvió como diagnóstico y
motivó la incorporación de las cotas. La segunda, con los modelos con cotas,
es la que alimenta el portafolio. En ambas campañas se emplearon tres
\textit{solvers} a través de MiniZinc~2.9.1, cuyas versiones y técnicas
principales se indican en la Tabla~\ref{tab:solvers}.

\begin{table}[H]
  \centering
  \small
  \singlespacing
  \caption{\textit{Solvers} utilizados.}
  \label{tab:solvers}
  \begin{tabularx}{\textwidth}{l c >{\raggedright\arraybackslash}X}
    \toprule
    \textbf{\textit{Solver}} & \textbf{Versión} & \textbf{Técnica principal} \\
    \midrule
    Chuffed & 0.13.2 & Generación perezosa de cláusulas (\textit{Lazy Clause Generation}) \\
    Gecode  & 6.3.0  & Propagación y búsqueda sobre dominios finitos \\
    OR-Tools CP-SAT & 9.12.4544 & Híbrido: generación perezosa de cláusulas sobre un núcleo SAT, combinada con técnicas de programación lineal \\
    \bottomrule
  \end{tabularx}
\end{table}

\subsection{Protocolo común}

Todas las ejecuciones siguieron el mismo protocolo. MiniZinc impone el límite
de 300 segundos con la opción \texttt{--time-limit}, que incluye la traducción a
FlatZinc y la búsqueda, y el programa de ejecución termina el árbol de procesos
si se supera ese límite en más de 20 segundos. Se usó la opción \texttt{-a},
para que el \textit{solver} reporte cada solución mejorada (sin ella, una
ejecución que agota el tiempo no reporta ninguna), y \texttt{-f}, que permite a
cada \textit{solver} usar su propia estrategia, salvo en el modelo con
estrategia de búsqueda. Ningún \textit{solver} se ejecutó con varios hilos, lo
que se verificó con la razón entre el tiempo de CPU y el de pared en la
instancia \texttt{dmu01} (0.92 para Chuffed, 0.93 para Gecode y 0.92 para
CP-SAT). Para cada par (instancia, configuración) se registra el último
\textit{makespan} reportado, el tiempo de pared, si se probó la optimalidad y
el estado final (\texttt{OPTIMAL}, \texttt{SATISFIABLE} si se encontró al
menos una solución sin probar que es óptima, o \texttt{UNKNOWN} si no se
encontró ninguna).

\subsection{Campaña con el modelo original}

La primera campaña ejecutó Chuffed, Gecode y CP-SAT con el modelo original
sobre las 242 instancias. Los resultados globales de esta exploración inicial
se sintetizan en la Tabla~\ref{tab:desempeno_300}, donde el \textit{gap}
corresponde a la definición de la Sección~\ref{sec:bks_gap}.

\begin{table}[H]
  \centering
  \small
  \singlespacing
  \caption{Desempeño de los \textit{solvers} con el modelo original
  (242 instancias, 300 segundos).}
  \label{tab:desempeno_300}
  \begin{tabular}{l c c c}
    \toprule
    & \textbf{Chuffed} & \textbf{Gecode} & \textbf{CP-SAT} \\
    \midrule
    Instancias con solución             & 242     & 242    & 242 \\
    Óptimos probados                     & 32      & 7      & 69 \\
    Instancias con \textit{gap} de 0\,\% & 32      & 7      & 76 \\
    \textit{Gap} medio (\%)              & 637.06  & 25.03  & 4.22 \\
    \textit{Gap} mediano (\%)            & 580.30  & 26.06  & 2.79 \\
    Tiempo medio (s)                     & 269.80  & 293.08 & 228.31 \\
    \bottomrule
  \end{tabular}
  \par\vspace{2pt}{\footnotesize Sin las seis instancias con $lb = ub$, el
  \textit{gap} medio es de 648.3, 25.6 y 4.3\,\% (Tabla~\ref{tab:efecto_cotas}).}
\end{table}

CP-SAT obtiene un \textit{makespan} igual o mejor que el de Chuffed y Gecode en
las 242 instancias, y estrictamente mejor en 210. Su ventaja generalizada sobre
ambos hace innecesaria una regla de selección, por lo que se dejó como
referencia externa (Sección~\ref{sec:composicion}). Chuffed prueba más óptimos
que Gecode en las instancias pequeñas, pero en 199 instancias su \textit{gap}
supera el 100\,\%, dado que con el horizonte trivial acepta como primeras
soluciones programas casi secuenciales y no logra mejorarlos dentro del límite.
Este comportamiento obedece a la amplitud del dominio en la formulación inicial
y motivó la incorporación de las cotas (Sección~\ref{sec:cotas}). Con ese
modelo el portafolio tenía también poco margen, ya que Gecode obtiene el mejor
\textit{makespan} entre los dos en 206 instancias, Chuffed en 29 y empatan en 7,
y la ventaja de Chuffed se limita a instancias pequeñas en las que ambos están
cerca del óptimo.
Estos \textit{solvers}, con el modelo original y su búsqueda por defecto, son
el punto de comparación de la propuesta, y la Tabla~\ref{tab:portafolios} los
contrasta con el portafolio construido sobre el modelo con cotas.

\subsection{Campaña con los modelos con cotas}

La segunda campaña ejecutó las tres configuraciones del portafolio
(Sección~\ref{sec:composicion}), integradas por Chuffed y Gecode con el modelo
con cotas y por Chuffed con el modelo con estrategia de búsqueda (Chuffed V2).
Se usó el programa de ejecución \texttt{src/05\_ejecutor}, diseñado para que los
resultados fueran verificables y la campaña pudiera repetirse. Antes de
empezar, comprueba las versiones de MiniZinc y de los \textit{solvers}
(Tabla~\ref{tab:solvers}) y resuelve \texttt{ft06}, cuyo óptimo (55) es
conocido. Con la opción \texttt{--output-time} de MiniZinc registra la
\textit{trayectoria} de cada ejecución, es decir, el instante en que se
encontró cada solución mejorada, que permite conocer $C(a, i, \tau)$ para
cualquier $\tau$ y evaluar cualquier \textit{schedule} sin volver a ejecutar
los \textit{solvers} (Sección~\ref{sec:formulacion_schedule}). También
comprueba que cada \textit{makespan} esté entre $lb$ y $ub$, marca las
ejecuciones sin solución y guarda cada resultado en cuanto termina, de modo
que una campaña interrumpida continúa donde quedó.

Las 708 ejecuciones (236 instancias por tres configuraciones) terminaron sin
errores. Ningún \textit{makespan} quedó fuera del intervalo $[lb, ub]$, y
todos los óptimos declarados coinciden con la mejor solución conocida. Una
ejecución sin solución significa que la configuración no encontró, en 300
segundos, un programa al menos tan bueno como el heurístico, situación que
ocurrió en 12 instancias con Chuffed, todas de 750 operaciones o más, en 24 con
Chuffed V2 y en 65 con Gecode, en ambos casos sobre todo en instancias grandes.
En esos casos el resultado de la configuración es la $ub$. Los resultados de
esta campaña se analizan en el Capítulo~\ref{cap:resultados} y se listan por
instancia en el Anexo~\ref{anx:resultados}.

Con el mismo programa y el mismo protocolo se ejecutó también CP-SAT con el
modelo con cotas, como \textit{solver} de referencia. Sus 236 ejecuciones
fueron válidas y encontraron solución en todas las instancias.

%------------------------------------------------------------------------------
\section{Preprocesamiento y base de conocimiento}
\label{sec:preprocesamiento}
%------------------------------------------------------------------------------

El preprocesamiento cruza por nombre de instancia las cotas, las mejores
soluciones conocidas, las características y los resultados de la campaña con
cotas. Las 242 instancias se vincularon en todas ellas.

\subsection{Mejores soluciones conocidas y brecha de optimalidad}
\label{sec:bks_gap}

Para medir la calidad de una solución se compara su \textit{makespan} con la
mejor solución conocida (BKS, por \textit{Best Known Solution}) de la
literatura, tomada del repositorio de \citet{weise2019jssp}, que recopila los
mejores resultados publicados hasta 2020. Para cada instancia $i$ y
configuración $a$, la brecha de optimalidad o \textit{gap} es
\begin{equation}
  \mathit{gap}(a,i) = 100 \cdot \frac{C_{\max}(a,i) - \mathit{BKS}(i)}{\mathit{BKS}(i)},
  \label{eq:gap}
\end{equation}
expresado en porcentaje, donde $C_{\max}(a,i)$ es el \textit{makespan}
obtenido por $a$ en $i$. Un \textit{gap} de 0\,\% indica que la configuración
alcanzó la mejor solución conocida. En las 156 instancias con óptimo
conocido, la BKS coincide con el óptimo, mientras que en las demás corresponde
a la mejor solución publicada, de modo que el \textit{gap} es una cota inferior
de la distancia real al óptimo.

\subsection{Limpieza y normalización de características}
\label{sec:limpieza}

\texttt{mzn2feat} reporta como \texttt{nan} o \texttt{inf} algunos
estadísticos que no se pueden calcular, por ejemplo cuando implican una
división por cero. Estos valores se trataron como faltantes. Solo los tenían
\texttt{c\_cv\_domdeg\_cons}, en 117 instancias, y \texttt{o\_dom\_std}, en
241, y se imputan con la mediana de la característica en el conjunto de
entrenamiento. Se eliminaron después las 30 características de
\texttt{mzn2feat} con varianza cero en las 242 instancias (anotaciones de
búsqueda, tipos de variables y restricciones que el modelo no usa y algunos
estadísticos que el modelo fija), con lo que quedaron 65 activas
(Tabla~\ref{tab:categorias_features}). Las 20 características propias no
tienen valores faltantes ni constantes.

Por último, siguiendo a SUNNY \citep{amadini2014sunny}, cada característica
$x$ se escala al intervalo $[-1, 1]$,
\begin{equation}
  x' = 2\,\frac{x - x_{\min}}{x_{\max} - x_{\min}} - 1,
  \label{eq:normalizacion}
\end{equation}
donde $x_{\min}$ y $x_{\max}$ son el mínimo y el máximo en el conjunto de
entrenamiento, y los valores de una instancia nueva que quedan fuera de ese
rango se recortan a $[-1, 1]$. Así se evita que las características con
valores grandes dominen la distancia euclidiana de los métodos de vecinos
más cercanos. La imputación y la normalización se calculan dentro de cada
partición de entrenamiento, para que la instancia evaluada no influya en su
propia transformación (Sección~\ref{sec:validacion}).

\subsection{Base de conocimiento}

Un aporte metodológico de este trabajo es la construcción de una base de
conocimiento con las trayectorias completas de soluciones de cada
configuración, lo que permite evaluar de forma exacta cualquier
\textit{schedule} por simulación sin repetir la ejecución de los
\textit{solvers}. Para cada una de las 242 instancias, la base reúne sus datos
generales (familia, $n$, $m$, número de operaciones, $lb$, $ub$, regla que
produjo la $ub$ y BKS), sus 85 características (65 de \texttt{mzn2feat} y 20
propias) y, para cada configuración, el \textit{makespan} final, si probó el
óptimo y en qué tiempo, el estado de la ejecución y la trayectoria completa
de soluciones. A diferencia de la base de conocimiento de SUNNY, que guarda
solo el resultado final de cada \textit{solver}, las trayectorias indican
cuánto tiempo necesitó cada configuración para alcanzar cada nivel de calidad,
que es justamente la información que requiere un \textit{schedule}. La base de
conocimiento se construye con el programa \texttt{src/06\_analisis\_campana.py},
que además produce los análisis descriptivos del
Capítulo~\ref{cap:resultados}. Con esto concluye la construcción del conjunto
de datos de entrenamiento (actividades A8 a A10).

%==============================================================================
\chapter{Diseño e implementación del portafolio}
\label{cap:portafolio}
%==============================================================================
% Objetivos específicos 1 y 3. Actividades A1 a A4 y A11 a A15.

Con la base de conocimiento disponible, el problema pasa a ser cómo elegir y
ordenar las configuraciones ante una instancia nueva dentro del presupuesto de
$T = 300$ segundos. Siguiendo la formulación de la
Sección~\ref{sec:formulacion_schedule}, el sistema integra en un mismo flujo la
preparación de la instancia, la extracción de sus características y la
construcción del \textit{schedule} $\sigma(i)$, ya sea mediante reglas de
referencia, métodos de vecinos más cercanos o modelos de aprendizaje
supervisado.

%------------------------------------------------------------------------------
\section{Arquitectura del sistema}
\label{sec:arquitectura}
%------------------------------------------------------------------------------

\subsection{Flujo de una ejecución}

El recorrido de una instancia nueva por el portafolio, desde su recepción hasta
la entrega de la mejor solución, sigue las seis etapas esquematizadas en la
Figura~\ref{fig:flujo_portafolio}:

\begin{enumerate}
  \item En la etapa de preparación, la instancia se convierte al formato
    \texttt{.dzn} (Sección~\ref{sec:instancias}) y se calculan sus cotas $lb$
    y $ub$ (Sección~\ref{sec:cotas}), dejando guardado el programa heurístico
    que produce la $ub$ como solución inicial.
  \item Si en esa verificación inicial se cumple que $lb = ub$, el programa
    heurístico ya es óptimo y el portafolio termina sin ejecutar ningún
    \textit{solver}.
  \item En caso contrario, se calcula el vector de características de la
    instancia y se escala con los mínimos y máximos de la base de
    conocimiento.
  \item A partir de ese vector, una estrategia de selección construye el
    \textit{schedule} definiendo qué configuraciones ejecutar, en qué orden y
    durante cuánto tiempo; en la estrategia principal la decisión se toma con
    tres modelos de aprendizaje automático
    (Sección~\ref{sec:modelos_prediccion}).
  \item A continuación, las configuraciones seleccionadas se ejecutan una tras
    otra con el modelo con cotas, deteniendo el proceso si alguna prueba el
    óptimo o pasando a la siguiente cuando agota su tiempo.
  \item Al concluir la secuencia, se entrega la mejor solución encontrada por
    cualquiera de las configuraciones o, si ninguna mejoró la heurística, el
    programa heurístico.
\end{enumerate}

\begin{figure}[H]
  \centering
  \includegraphics[height=0.55\textheight]{flujo_portafolio.pdf}
  \caption{Flujo de ejecución del portafolio para una instancia nueva.}
  \label{fig:flujo_portafolio}
\end{figure}

\subsection{Componentes}

El sistema se organiza en los componentes de la Tabla~\ref{tab:componentes}.
Cada uno es un programa independiente del repositorio del trabajo
(Anexo~\ref{anx:reproduccion}), de modo que cualquier etapa puede repetirse
por separado. El ejecutor es el componente que construyó la base de
conocimiento (Sección~\ref{sec:ejecucion_solvers}), el cual registra, además
del resultado final, la trayectoria completa de soluciones de cada ejecución
para evaluar cualquier \textit{schedule} con la
ecuación~\eqref{eq:resultado_schedule}.

\begin{table}[H]
  \centering
  \small
  \singlespacing
  \caption{Componentes del sistema y programas que los implementan.}
  \label{tab:componentes}
  \begin{tabularx}{\textwidth}{>{\raggedright\arraybackslash}p{3.4cm} >{\raggedright\arraybackslash}X l}
    \toprule
    \textbf{Componente} & \textbf{Función} & \textbf{Programa (\texttt{src/})} \\
    \midrule
    Conversión de instancias & Traduce el formato de la OR-Library a
      \texttt{.dzn}, con y sin cotas & \texttt{01\_convertir\_instancias.py} \\
    Cálculo de cotas & Algoritmo de Giffler y Thompson con seis reglas y
      verificación de cada programa & \texttt{02\_generar\_cotas.py} \\
    Características genéricas & Extracción con \texttt{mzn2feat} &
      \texttt{03\_caracteristicas\_mzn2feat.py} \\
    Características propias & Características específicas del JSSP &
      \texttt{04\_caracteristicas\_propias.py} \\
    Ejecutor & Lanza cada configuración con MiniZinc, aplica el límite de
      tiempo y registra las trayectorias & \texttt{05\_ejecutor/} \\
    Estrategias de selección & SBS, regla por tamaño, \textit{Winner Takes
      All}, SUNNY y \textit{schedule} estático & \texttt{07\_portafolios.py} \\
    \textit{Schedule} con aprendizaje & Tres modelos de predicción y
      construcción del \textit{schedule} & \texttt{08\_modelo\_ml.py} \\
    \bottomrule
  \end{tabularx}
\end{table}

Para medir el aporte de esta propuesta, el portafolio se compara también con
los mismos \textit{solvers} ejecutados con el modelo original, sin cotas
(Sección~\ref{sec:ejecucion_solvers} y Tabla~\ref{tab:portafolios}).

%------------------------------------------------------------------------------
\section{Composición del portafolio}
\label{sec:composicion}
%------------------------------------------------------------------------------

La primera campaña de ejecución, con el modelo original, incluyó a Chuffed,
Gecode y OR-Tools CP-SAT (Sección~\ref{sec:ejecucion_solvers}). En ella CP-SAT
obtuvo la mejor solución en las 242 instancias, y en 210 de ellas fue el único
en alcanzarla. Dado que esa ventaja se mantiene en casi todo el conjunto, para
estudiar la selección entre alternativas complementarias y con la aprobación de
la dirección del trabajo, el portafolio activo quedó compuesto por tres
configuraciones basadas en Chuffed y Gecode (Tabla~\ref{tab:configuraciones}).

\begin{table}[H]
  \centering
  \small
  \singlespacing
  \caption{Configuraciones del portafolio.}
  \label{tab:configuraciones}
  \begin{tabularx}{\textwidth}{l l >{\raggedright\arraybackslash}X}
    \toprule
    \textbf{Configuración} & \textbf{Modelo} & \textbf{Fortaleza} \\
    \midrule
    Chuffed & Con cotas & Prueba la optimalidad de las instancias pequeñas y
      medianas \\
    Chuffed V2 & Con cotas y estrategia de búsqueda & Encuentra mejores
      soluciones en las instancias grandes \\
    Gecode & Con cotas & Encuentra buenas soluciones en los primeros segundos
      en algunas instancias grandes \\
    \bottomrule
  \end{tabularx}
\end{table}

Chuffed combina búsqueda con aprendizaje de cláusulas y, con su búsqueda
libre, prueba el óptimo de muchas instancias pequeñas y medianas. Chuffed V2 es
el mismo \textit{solver} con la estrategia de búsqueda de la
Sección~\ref{sec:modelo_busqueda}, que prueba menos óptimos pero obtiene
mejores soluciones en las instancias grandes. Gecode, un \textit{solver} de
propagación clásica sin aprendizaje, es el más débil de los tres, pero aporta
diversidad, ya que es un \textit{solver} distinto y en algunas instancias
grandes encuentra la mejor solución en los primeros segundos.
Incluir dos configuraciones de un mismo \textit{solver} es una práctica
habitual en los portafolios basados en configuraciones, como Hydra
\citep{xu2010hydra}, cuyo criterio principal es la complementariedad entre sus
miembros. En la campaña con el modelo con cotas, las tres configuraciones
muestran desempeños complementarios, con ventaja de Chuffed en las instancias
pequeñas y medianas y de Chuffed V2 en las grandes
(Capítulo~\ref{cap:resultados}). CP-SAT se conserva como \textit{solver} de
referencia externo, para situar el desempeño del portafolio respecto a un
\textit{solver} del estado del arte.

%------------------------------------------------------------------------------
\section{Estrategias de selección}
\label{sec:estrategias}
%------------------------------------------------------------------------------

Se implementaron cinco estrategias, de menor a mayor complejidad. Todas usan
únicamente información de la base de conocimiento, que comprende para cada
configuración e instancia de entrenamiento el \textit{makespan} obtenido en
función del tiempo, $C(a, i, \tau)$, y el registro de si probó el óptimo.

\subsection{Mejor \textit{solver} individual}
La estrategia más simple ejecuta siempre la configuración con menor brecha
media en el conjunto de entrenamiento (SBS). En todas las particiones de la
validación resultó ser Chuffed V2. Es la referencia que cualquier otra estrategia
debe superar.

\subsection{Regla por tamaño}
La segunda estrategia usa una sola característica, el número de operaciones
$n \cdot m$. Busca en el entrenamiento el par de configuraciones $(a, b)$ y el
umbral $u$ que minimizan la brecha media de la regla ``ejecutar $b$ si
$n \cdot m \ge u$ y $a$ en otro caso'', y la aplica a la instancia nueva.
Sirve para medir cuánto aporta un modelo con muchas características frente a
una regla que un experto podría escribir a mano.

\subsection{\textit{Winner Takes All} con vecinos más cercanos}
\label{sec:wta}
La estrategia \textit{Winner Takes All} elige una sola configuración y le
asigna los 300 segundos. La elección se hace con el método de los $k$
vecinos más cercanos \citep{cover1967}, escalando las características a
$[-1, 1]$ según los mínimos y máximos del entrenamiento para buscar las $k$
instancias de entrenamiento más cercanas a la nueva según la distancia
euclidiana y elegir la configuración con menor brecha media en ellas. Se usa
$k = 15$, la raíz cuadrada del tamaño del conjunto de entrenamiento
($\sqrt{235} \approx 15$), que es el criterio con el que SUNNY-CP fija el
tamaño del vecindario \citep{amadini2015sunnycp}, y su sensibilidad se
analiza en el Capítulo~\ref{cap:resultados}.

\subsection{\textit{Schedule} secuencial basado en SUNNY}
\label{sec:sunny}
SUNNY \citep{amadini2014sunny, amadini2015sunnycp} construye un
\textit{schedule} para cada instancia a partir de sus vecinos más cercanos.
Su versión original está pensada para problemas de satisfacción, donde una
instancia solo puede quedar resuelta o sin resolver. Como aporte de este
trabajo se implementó una adaptación de SUNNY al JSSP como problema de
optimización, estructurada en cinco pasos:

\begin{enumerate}
  \item \textbf{Vecindario.} Se obtienen los $k = 15$ vecinos más cercanos,
    igual que en \textit{Winner Takes All}.
  \item \textbf{Cobertura.} Se dice que una configuración \textit{cubre} a un
    vecino si alcanza en él el mejor \textit{makespan} del portafolio. Se
    elige el subconjunto de configuraciones que cubre más vecinos y, en caso de
    empate, el de menos configuraciones, seguido del que necesita menos tiempo
    para alcanzar esos valores.
  \item \textbf{Asignación de tiempo.} Cada configuración del subconjunto
    recibe una porción del presupuesto proporcional al número de vecinos que
    cubre.
  \item \textbf{Respaldo.} Los vecinos que ninguna configuración cubre se
    asignan a la configuración con menor brecha media en el vecindario, que
    recibe su porción adicional de tiempo. Este paso proviene del SUNNY
    original; como aquí la cobertura se mide respecto al mejor valor del
    propio portafolio, todo vecino queda cubierto por alguna configuración y
    el paso no llega a activarse, aunque se conserva por completitud.
  \item \textbf{Orden.} Las configuraciones se ejecutan de menor a mayor
    tiempo medio necesario para alcanzar el mejor valor en los vecinos que
    cubren.
\end{enumerate}

\subsection{\textit{Schedule} estático}
\label{sec:schedule_estatico}
Otro aporte realizado en este trabajo es el diseño de un \textit{schedule}
estático con un tramo inicial corto de Gecode, el cual aplica a todas las
instancias el reparto fijo del tiempo con mejor desempeño promedio en el
entrenamiento. Los candidatos evaluados abarcaron todos los pares de
configuraciones en cualquier orden, $[(a, x), (b, T - x)]$, con
$x \in \{2, 5, 10, 15, 20, 30, 45, 60, 90, 120, 150, 200\}$ segundos, los
casos extremos $x = 0$ y $x = 300$ para Gecode y Chuffed, que equivalen a
ejecutar una sola configuración, y tríos que empiezan con 2 segundos de
Gecode, $[(\text{Gecode}, 2), (a, y), (b, T - 2 - y)]$, con
$y \in \{30, 60, 90, 120, 150, 180\}$. Los tríos se incluyeron porque Gecode,
cuando supera a las otras configuraciones, suele hacerlo con las soluciones
que encuentra en los primeros segundos. En todas las particiones se eligió el
mismo \textit{schedule}:
\[
  \sigma_{\mathrm{est}} = \bigl[ (\text{Gecode}, 2), (\text{Chuffed V2}, 60),
  (\text{Chuffed}, 238) \bigr].
\]
Estos valores provienen de una calibración empírica: los tiempos candidatos
se fijaron a mano, cada combinación se evaluó por simulación con las
trayectorias de la base de conocimiento, sin nuevas ejecuciones, y se eligió
la de menor \textit{gap} medio en las instancias de entrenamiento.
Esta estrategia sirve de referencia para medir cuánto aportan las
características de la instancia frente a un reparto fijo bien calibrado.

%------------------------------------------------------------------------------
\section{Modelos de predicción y \textit{schedule} con aprendizaje automático}
\label{sec:modelos_prediccion}
%------------------------------------------------------------------------------

Un aporte central de este trabajo es la estrategia principal del portafolio,
que construye el \textit{schedule} de cada instancia combinando las
predicciones de tres modelos de aprendizaje supervisado. En la literatura de
selección de algoritmos es habitual predecir el desempeño de cada algoritmo a
partir de las características de la instancia \citep{xu2008satzilla,
hutter2014}. Aquí cada modelo responde una de las preguntas que plantea el
\textit{schedule}, determinando si conviene darle tiempo a una configuración
para probar el óptimo, cuánto tiempo necesita y cuál de las configuraciones
dará la mejor solución cuando ninguna lo pruebe.

\subsection{Modelos}
Los tres modelos son bosques aleatorios \citep{breiman2001} de 300 árboles,
con al menos dos ejemplos por hoja, implementados con scikit-learn
\citep{pedregosa2011}. Usan como entrada las características de la instancia
(las de \texttt{mzn2feat} y las propias, escaladas a $[-1, 1]$) y, en los dos
primeros, un indicador de la configuración, de modo que un mismo modelo
aprende el comportamiento de todas.

\begin{itemize}
  \item \textbf{Clasificador de resolución.} Estima la probabilidad
    $p_{\mathrm{res}}(a, i)$ de que la configuración $a$ pruebe el óptimo de
    $i$ dentro de $T$.
  \item \textbf{Regresor de tiempo.} Estima el tiempo $\hat{t}(a, i)$ que
    tardará $a$ en probar el óptimo. Se entrena solo con las ejecuciones que
    lo probaron y sobre el logaritmo del tiempo, porque los tiempos van de
    décimas de segundo a varios minutos. Restringir el entrenamiento a esas
    ejecuciones evita tratar como tiempos reales los valores censurados en el
    límite de 300 segundos, aunque, como advierten \citet{hutter2014},
    descartar los datos censurados tiende a subestimar los tiempos; por eso el
    \textit{schedule} asigna a cada configuración el doble del tiempo predicho.
  \item \textbf{Clasificador de la mejor configuración.} Estima, para cada
    configuración $a$, la probabilidad $p_{\mathrm{mej}}(a, i)$ de que sea la
    que obtiene el menor \textit{makespan} al cabo de los 300 segundos. Si
    varias empatan, se toma como mejor la que alcanza ese valor primero.
\end{itemize}

\subsection{Construcción del \textit{schedule}}
Con las tres predicciones, el \textit{schedule} de una instancia se construye
en dos pasos:

\begin{enumerate}
  \item Primero se ejecutan, de menor a mayor tiempo predicho, las
    configuraciones que probablemente prueban el óptimo, definidas por
    $p_{\mathrm{res}}(a, i) \ge 0.5$. Cada una recibe el doble de su tiempo
    predicho, $2\,\hat{t}(a, i)$, como margen ante el error del regresor, sin
    exceder el tiempo disponible.
  \item A continuación, el tiempo restante $R$ se reparte según la probabilidad
    de que cada configuración sea la mejor. Cada una recibe
    $\max\bigl(m_a,\ p_{\mathrm{mej}}(a, i)\,R\bigr)$ segundos, donde $m_a$ es
    el tiempo que le da el \textit{schedule} estático cuando ese tramo es corto
    (menos de 150 segundos), equivalente a 2 segundos para Gecode y 60 para
    Chuffed V2. Las asignaciones se reescalan para que sumen $R$ y las
    configuraciones se ejecutan de menor a mayor tiempo asignado.
\end{enumerate}

\noindent Así, en una instancia pequeña que Chuffed prueba en pocos segundos,
el \textit{schedule} empieza por Chuffed con un tiempo corto, mientras que en
una instancia grande en la que se espera que gane Chuffed V2, esta
configuración recibe la mayor parte del presupuesto y, cuando el modelo carece
de una predicción clara, el reparto se aproxima al estático.

Todas las estrategias se evalúan con validación \textit{leave-one-out}, de modo
que los modelos, los vecinos, el umbral de la regla por tamaño y el
\textit{schedule} estático se calculan sin la instancia evaluada
(Capítulo~\ref{cap:diseno}).

%------------------------------------------------------------------------------
\section{Entorno de desarrollo y verificación}
\label{sec:entorno}
%------------------------------------------------------------------------------

El sistema se implementó en Python con pandas, NumPy y scikit-learn
\citep{pedregosa2011}, y usa MiniZinc 2.9.1 con Chuffed 0.13.2, Gecode 6.3.0
y OR-Tools CP-SAT 9.12. Para que los resultados fueran confiables, cada
programa construido por la heurística se comprueba contra las restricciones de
precedencia y de capacidad, y se verifica que $lb \le ub$ y que $lb$ no supere
la mejor solución conocida en ninguna instancia. Ningún \textit{makespan}
registrado es menor que la mejor solución conocida de su instancia
\citep{weise2019jssp}, y todo óptimo declarado por un \textit{solver} coincide
con ella. Los programas de cotas y de características propias reproducen
exactamente los archivos usados en el análisis, y todas las tablas y figuras
de resultados se generan con programas del repositorio
(Anexo~\ref{anx:reproduccion}).

%==============================================================================
\chapter{Diseño experimental}
\label{cap:diseno}
%==============================================================================
% Objetivos específicos 1 y 4.

Para comprobar si el reparto del tiempo entre configuraciones y el uso de las
características de cada instancia mejoran la resolución del JSSP frente a
ejecutar siempre un único \textit{solver}, se diseñó una evaluación controlada
sobre las instancias de referencia. Antes de pasar al análisis cuantitativo del
Capítulo~\ref{cap:resultados}, aquí se precisan las preguntas experimentales,
el esquema de validación cruzada, los métodos comparados y las métricas
empleadas.

%------------------------------------------------------------------------------
\section{Preguntas de evaluación}
\label{sec:preguntas}
%------------------------------------------------------------------------------

La evaluación se organizó en torno a cinco preguntas, derivadas de los
objetivos específicos:

\begin{enumerate}[label=\textbf{P\arabic*.}, leftmargin=*]
  \item \textbf{Efecto de las cotas.} ¿Cuánto mejora el desempeño de los
    \textit{solvers} al pasar del modelo original al modelo con cotas?
  \item \textbf{Complementariedad.} ¿Las configuraciones del portafolio son
    complementarias? ¿En qué tipo de instancias es mejor cada una y cuánto
    margen deja eso para un portafolio?
  \item \textbf{\textit{Winner Takes All} frente a \textit{schedule}
    secuencial} (objetivo específico~1). ¿Repartir el tiempo entre varias
    configuraciones da mejores resultados que elegir una sola?
  \item \textbf{Aporte de las características y del aprendizaje automático}
    (objetivo específico~3). ¿Las estrategias que usan las características
    de la instancia superan a las que no las usan, como el mejor
    \textit{solver} individual, la regla por tamaño o el \textit{schedule}
    estático?
  \item \textbf{Comparación con el estado del arte} (objetivo
    específico~4). ¿Cómo se sitúa el portafolio frente a CP-SAT, un
    \textit{solver} de referencia que no forma parte de él?
\end{enumerate}

%------------------------------------------------------------------------------
\section{Instancias y esquema de validación}
\label{sec:validacion}
%------------------------------------------------------------------------------

La evaluación usa las 236 instancias en las que $lb < ub$. Las seis restantes
quedan resueltas por las cotas (Sección~\ref{sec:cotas_instancias}), puesto que
todos los métodos obtienen en ellas el óptimo sin ejecutar ningún
\textit{solver}, así que incluirlas no cambiaría la comparación.

Como el conjunto de instancias es pequeño, se usó validación cruzada dejando
una instancia fuera (\textit{leave-one-out}). Para cada una de las 236
instancias se toman las otras 235 como entrenamiento, se calcula con ellas
todo lo que el método necesita aprender (la imputación y la normalización de
las características, los vecinos más cercanos, el umbral de la regla por
tamaño, el \textit{schedule} estático y los modelos de aprendizaje automático)
y se construye el \textit{schedule} de la instancia excluida, que se evalúa con
sus resultados reales. Así, ninguna decisión sobre una instancia usa
información de esa misma instancia, y las comparaciones entre métodos son
pareadas, porque todos se evalúan sobre las mismas 236 instancias.

\subsection{Evaluación de \textit{schedules} por simulación}

Cada \textit{schedule} se evaluó directamente sobre las trayectorias
registradas en la campaña con cotas, mediante la
ecuación~\eqref{eq:resultado_schedule}. Para una configuración $a$ con un
tiempo asignado $\tau$, su resultado es el mejor \textit{makespan} que $a$
había encontrado en el instante $\tau$ de su ejecución de 300 segundos, o la
$ub$ si todavía no había encontrado ninguno. Si $a$ probó el óptimo antes de
$\tau$, el \textit{schedule} termina en ese momento. Esta simulación es exacta
porque las configuraciones de un \textit{schedule} se ejecutan de forma
independiente, por lo cual la ejecución de $a$ durante $\tau$ segundos dentro
del \textit{schedule} coincide con el tramo inicial de su ejecución de 300
segundos.

La simulación permite comparar muchas estrategias y variantes con el mismo
costo computacional que una sola campaña, y garantiza que las diferencias
entre métodos se deban a sus decisiones y no a variaciones entre ejecuciones.

%------------------------------------------------------------------------------
\section{Enfoques comparados y referencias}
\label{sec:enfoques}
%------------------------------------------------------------------------------

En la Tabla~\ref{tab:metodos} se clasifican los métodos evaluados, cuya
implementación fue descrita en el Capítulo~\ref{cap:portafolio}. Se agrupan en
referencias, estrategias que no usan características, estrategias de selección
de una configuración (\textit{Winner Takes All}) y estrategias que construyen
un \textit{schedule}.

WTA $k$-NN y SUNNY se evaluaron con los tres conjuntos de características de
la Sección~\ref{sec:features}, formados por las de \texttt{mzn2feat}, las
propias y la unión de ambas. El \textit{schedule} con ML usa la unión. El
número de vecinos es $k = 15$, y la sensibilidad a este valor se evaluó con
$k \in \{3, 5, 7, 10, 15, 20, 30\}$.

El VBS no es un método aplicable, porque requiere conocer de antemano el
resultado de cada configuración, de modo que representa la cota superior de lo
que puede lograr cualquier estrategia con las mismas configuraciones. CP-SAT,
en cambio, es un \textit{solver} real, pero no forma parte del portafolio
(Sección~\ref{sec:composicion}) y su papel consiste en situar los resultados
respecto al estado del arte.

\begin{table}[H]
  \centering
  \small
  \singlespacing
  \caption{Métodos comparados.}
  \label{tab:metodos}
  \begin{tabularx}{\textwidth}{l >{\raggedright\arraybackslash}X c c}
    \toprule
    \textbf{Método} & \textbf{Descripción} & \textbf{Usa caract.} & \textbf{Varias config.} \\
    \midrule
    \multicolumn{4}{l}{\textit{Referencias}} \\
    Configuraciones individuales & Chuffed, Gecode y Chuffed V2 durante 300\,s & No & No \\
    SBS & La configuración con menor \textit{gap} medio en el entrenamiento & No & No \\
    VBS & La mejor configuración en cada instancia (oráculo) & -- & No \\
    CP-SAT & \textit{Solver} de referencia externo, fuera del portafolio & No & No \\
    \midrule
    \multicolumn{4}{l}{\textit{Sin características}} \\
    Regla por tamaño & Una configuración para instancias pequeñas y otra para grandes & Solo $n \cdot m$ & No \\
    \textit{Schedule} estático & El mismo \textit{schedule} para todas las instancias & No & Sí \\
    \midrule
    \multicolumn{4}{l}{\textit{Winner Takes All}} \\
    WTA $k$-NN & La configuración con menor \textit{gap} medio en los $k$ vecinos & Sí & No \\
    \midrule
    \multicolumn{4}{l}{\textit{Schedule secuencial}} \\
    SUNNY & \textit{Schedule} a partir de los $k$ vecinos más cercanos & Sí & Sí \\
    \textit{Schedule} con ML & \textit{Schedule} a partir de tres modelos de predicción & Sí & Sí \\
    \bottomrule
  \end{tabularx}
\end{table}

%------------------------------------------------------------------------------
\section{Métricas}
\label{sec:metricas}
%------------------------------------------------------------------------------

Se usaron las métricas de la Sección~\ref{sec:metricas_portafolios}, con
$T = 300$ segundos. La principal es el \textbf{\textit{gap} medio}
(ecuación~\eqref{eq:gap}), porque solo 36 de las 236 instancias se pueden
probar óptimas dentro del límite con alguna configuración y, en la gran
mayoría, lo relevante es la calidad de la mejor solución encontrada. Se
reportan también los \textbf{óptimos probados} y el \textbf{PAR10}
\citep{lindauer2019}, con
$10 \cdot T = 3000$ segundos para las instancias en que no se prueba el óptimo.
En un \textit{schedule}, el tiempo incluye el de las configuraciones
ejecutadas antes de la que lo prueba. La \textbf{brecha cerrada} es la
fracción del margen entre el SBS y el VBS que el método logra cerrar
\citep{kerschke2019, lindauer2019},
\begin{equation}
  \mathit{brecha} = 100 \cdot
  \frac{\overline{\mathit{gap}}_{\mathrm{SBS}} - \overline{\mathit{gap}}}
       {\overline{\mathit{gap}}_{\mathrm{SBS}} - \overline{\mathit{gap}}_{\mathrm{VBS}}},
  \label{eq:brecha}
\end{equation}
que vale 0\,\% para el SBS, 100\,\% para el VBS y es negativa si el método es
peor que usar siempre la mejor configuración individual. Por último, el número
de instancias en que el método obtiene un \textit{makespan} menor o mayor que
el SBS distingue a un método que mejora mucho en pocas instancias de uno que
mejora poco en muchas.

Para comparar dos métodos se usan sus diferencias de \textit{gap} instancia
por instancia. Se reportan la diferencia media con un intervalo de confianza
del 95\,\% obtenido por \textit{bootstrap} (10\,000 remuestreos) y el
$p$-valor de la prueba de rangos con signo de Wilcoxon sobre las instancias
en que los resultados difieren. El análisis descriptivo de la campaña incluye
asimismo el \textit{gap} medio de cada configuración en función del tiempo (10,
30, 60, 120 y 300 segundos), que muestra su comportamiento \textit{anytime}.

%------------------------------------------------------------------------------
\section{Protocolo de ejecución}
\label{sec:protocolo}
%------------------------------------------------------------------------------

Los datos provienen de la campaña con cotas de la
Sección~\ref{sec:ejecucion_solvers}, en la cual cada configuración se ejecutó
una vez sobre cada instancia, con un límite de 300 segundos y un solo núcleo.
Las ejecuciones se hicieron de dos en dos, cada una fijada a un núcleo físico
distinto del procesador. Las versiones del software y las características del
equipo se detallan en el Anexo~\ref{anx:configuracion}.

Los métodos de selección y los modelos de aprendizaje se implementaron en
Python con scikit-learn \citep{pedregosa2011}. Los bosques aleatorios usan una
semilla fija, y los empates en los vecinos más cercanos y en la elección de
configuraciones se resuelven con un orden fijo, de modo que todos los
resultados son deterministas.

%------------------------------------------------------------------------------
\section{Reproducibilidad}
\label{sec:reproducibilidad}
%------------------------------------------------------------------------------

Todo el proceso puede repetirse con el repositorio que acompaña este
documento, cuyos programas siguen el orden del trabajo. Los scripts \texttt{01}
a \texttt{04} convierten las instancias y calculan las cotas y las
características (Capítulo~\ref{cap:caracterizacion}), \texttt{05} ejecuta las
campañas, \texttt{06} construye la base de conocimiento y el análisis
descriptivo, \texttt{07} y \texttt{08} implementan las estrategias de selección
y los modelos de aprendizaje con su validación, y \texttt{09} genera las tablas
y figuras del Capítulo~\ref{cap:resultados}. Como los resultados de las
campañas se incluyen en el repositorio, los programas \texttt{06} a \texttt{09}
pueden ejecutarse sin repetirlas, y la mayor parte de su tiempo de cómputo
corresponde al entrenamiento de los modelos en cada partición de la
validación. La guía completa está en el Anexo~\ref{anx:reproduccion}.

%==============================================================================
\chapter{Resultados y análisis}
\label{cap:resultados}
%==============================================================================
% Objetivos específicos 1 y 4. Actividad A16.

La pregunta de investigación plantea determinar en qué medida un
\textit{schedule} basado en portafolios y guiado por aprendizaje automático
mejora la eficiencia y la robustez en el JSSP clásico. En este trabajo, esa
pregunta se responde para un portafolio concreto, formado por tres
configuraciones de Chuffed y Gecode (Chuffed, Gecode y Chuffed V2) con el
modelo con cotas, de modo que los resultados y las conclusiones se refieren a
ese portafolio y a las instancias estudiadas. Para responderla en el
orden de las cinco preguntas de la Sección~\ref{sec:preguntas}, se examinan a
continuación los resultados obtenidos sobre las 236 instancias con $lb < ub$
bajo un presupuesto de 300 segundos por instancia y validación
\textit{leave-one-out}, salvo que se indique otra cosa. Todas las tablas y
figuras se generan con el programa \texttt{src/09\_tablas\_figuras.py}.

%------------------------------------------------------------------------------
\section{Efecto de las cotas}
\label{sec:res_cotas}
%------------------------------------------------------------------------------

En la Tabla~\ref{tab:efecto_cotas} se contrasta el desempeño de Chuffed, Gecode
y CP-SAT entre el modelo original y el modelo con cotas.

\begin{table}[H]
  \centering
  \small
  \singlespacing
  \caption{Efecto de las cotas en el modelo (236 instancias, 300\,s).}
  \label{tab:efecto_cotas}
  \begin{tabular}{lrrrrrrrrr}
    \toprule
    \textit{Solver} & \multicolumn{2}{c}{Gap medio (\%)} & \multicolumn{3}{c}{Óptimos} & \multicolumn{2}{c}{$t$ mediano (s)} & Mejora & Empeora \\
    \cmidrule(lr){2-3}\cmidrule(lr){4-6}\cmidrule(lr){7-8}
     & Original & Cotas & Original & Cotas & Nuevos & Original & Cotas &  &  \\
    \midrule
    Chuffed & 648.3 & 16.8 & 28 & 36 & 8 & 60.1 & 2.8 & 207 & 1 \\
    Gecode & 25.6 & 21.7 & 5 & 6 & 1 & 2.6 & 0.4 & 172 & 51 \\
    CP-SAT (ref.) & 4.3 & 4.0 & 63 & 62 & 1 & 9.8 & 9.9 & 82 & 22 \\
    \bottomrule
  \end{tabular}
  \par\vspace{2pt}{\footnotesize $t$ mediano: tiempo hasta probar el óptimo en las instancias donde ambos modelos lo prueban. Mejora/Empeora: instancias con menor/mayor \textit{makespan} que con el modelo original.}
\end{table}


En Chuffed, la incorporación de las cotas modifica de manera sustancial el
desempeño. Su \textit{gap} medio pasa de 648.3\,\% a 16.8\,\%, puesto que con el
modelo original se quedaba en programas casi secuenciales en las instancias
grandes, mientras que con las cotas parte de soluciones al menos tan buenas como
la heurística. Prueba además la optimalidad en 8 instancias adicionales y, en
las instancias en que ambos modelos la prueban, el tiempo mediano baja de 60.1 a
2.8 segundos. Mejora en 207 de las 236 instancias y empeora solo en una.

En Gecode el efecto es más moderado. Su \textit{gap} medio baja de 25.6\,\% a
21.7\,\% y mejora en 172 instancias, pero empeora en 51. La razón obedece a la
segunda consecuencia de las cotas descrita en la
Sección~\ref{sec:modelo_cotas}, ya que Gecode no encuentra en 300 segundos un
programa mejor que el heurístico en 65 instancias, y en ellas su resultado es
la $ub$, a veces peor que la solución que encontraba con el modelo original.

En CP-SAT, que ya tenía un buen desempeño con el modelo original, la variación
es pequeña, con una reducción del \textit{gap} medio de 4.3\,\% a 4.0\,\%,
mejoras en 82 instancias, retrocesos en 22 y prácticamente los mismos óptimos
probados. CP-SAT calcula internamente sus propias cotas y soluciones iniciales,
de modo que la información adicional del modelo le aporta poco.

De este modo, la formulación con cotas sitúa a Chuffed como una configuración
competitiva dentro del portafolio y sirve de base para las evaluaciones
siguientes.

%------------------------------------------------------------------------------
\section{Desempeño de las configuraciones}
\label{sec:res_solvers}
%------------------------------------------------------------------------------

La Tabla~\ref{tab:individual} resume el desempeño de las tres configuraciones
del portafolio con un presupuesto de 300 segundos.

\begin{table}[H]
  \centering
  \small
  \singlespacing
  \caption{Desempeño de las configuraciones con el modelo con cotas (236 instancias, 300\,s).}
  \label{tab:individual}
  \begin{tabular}{lrrrrrr}
    \toprule
     & Chuffed & Gecode & Chuffed V2 & VBS & Heurística & CP-SAT (ref.) \\
    \midrule
    Óptimos probados & 36 & 6 & 18 & 36 & -- & 62 \\
    Gap medio (\%) & 16.76 & 21.68 & 15.08 & 12.83 & 26.96 & 3.97 \\
    PAR10 (s) & 2550 & 2924 & 2774 & 2550 & -- & 2224 \\
    Instancias donde gana & 100 & 3 & 100 & -- & -- & -- \\
    Sin solución & 12 & 65 & 24 & -- & -- & 0 \\
    \bottomrule
  \end{tabular}
  \par\vspace{2pt}{\footnotesize Se excluyen las 6 instancias con lb $=$ ub. Sin solución: el resultado es la ub. VBS: la mejor configuración del portafolio en cada instancia.}
\end{table}


Los resultados muestran un reparto claro de fortalezas entre las tres
configuraciones. Chuffed prueba la optimalidad en 36 instancias, el doble que
Chuffed V2, y encuentra solución en casi todas. Chuffed V2 tiene el menor
\textit{gap} medio (15.08\,\%) y constituye, por tanto, el SBS, al obtener la
mejor solución en 100 instancias, tantas como Chuffed, aunque prueba menos
óptimos. Gecode es la configuración más débil, ya que solo gana en 3 instancias
y no encuentra solución en 65. El VBS alcanza un \textit{gap} medio de
12.83\,\%, de modo que el margen disponible para el portafolio es de 2.24
puntos. La heurística sola tiene un \textit{gap} de 26.96\,\%.

\subsection{Complementariedad por tamaño y familia}

Al desagregar los resultados por número de operaciones
(Tabla~\ref{tab:por_tamano}) y cruzar el \textit{gap} de ambas versiones de
Chuffed (Figura~\ref{fig:dispersion}), se aprecia una complementariedad
estrechamente ligada al tamaño de la instancia.

\begin{table}[H]
  \centering
  \small
  \singlespacing
  \caption{Gap medio (\%) por tamaño de instancia.}
  \label{tab:por_tamano}
  \begin{tabular}{lrrrrrrrrr}
    \toprule
    Tamaño & Inst. & \multicolumn{5}{c}{Gap medio (\%)} & \multicolumn{3}{c}{Gana} \\
    \cmidrule(lr){3-7}\cmidrule(lr){8-10}
     &  & C & G & C2 & VBS & Margen & C & G & C2 \\
    \midrule
    6$\times$6 & 1 & 0.00 & 0.00 & 0.00 & 0.00 & 0.00 & 0 & 0 & 0 \\
    10$\times$5 & 5 & 0.00 & 4.39 & 1.91 & 0.00 & 0.00 & 0 & 0 & 0 \\
    15$\times$5 & 3 & 0.00 & 1.44 & 0.00 & 0.00 & 0.00 & 0 & 0 & 0 \\
    10$\times$10 & 18 & 0.07 & 9.04 & 3.15 & 0.07 & 0.00 & 9 & 0 & 0 \\
    20$\times$5 & 4 & 0.60 & 9.65 & 9.16 & 0.60 & 0.00 & 3 & 0 & 0 \\
    15$\times$10 & 5 & 0.21 & 15.39 & 2.85 & 0.21 & 0.00 & 4 & 0 & 0 \\
    20$\times$10 & 10 & 8.51 & 22.72 & 18.39 & 8.45 & 0.06 & 9 & 0 & 1 \\
    15$\times$15 & 15 & 2.10 & 16.08 & 6.01 & 2.10 & 0.00 & 15 & 0 & 0 \\
    30$\times$10 & 5 & 1.82 & 9.20 & 6.09 & 1.82 & 0.00 & 5 & 0 & 0 \\
    20$\times$15 & 28 & 12.29 & 24.23 & 16.63 & 11.86 & 0.43 & 22 & 0 & 6 \\
    20$\times$20 & 24 & 12.67 & 22.31 & 12.29 & 10.57 & 1.72 & 13 & 0 & 11 \\
    30$\times$15 & 20 & 24.91 & 28.57 & 21.26 & 19.50 & 1.76 & 8 & 0 & 12 \\
    50$\times$10 & 8 & 30.62 & 23.44 & 14.54 & 14.54 & 0.00 & 0 & 0 & 7 \\
    40$\times$15 & 10 & 32.06 & 29.46 & 23.02 & 22.89 & 0.12 & 2 & 0 & 8 \\
    30$\times$20 & 20 & 27.05 & 29.73 & 22.82 & 21.79 & 1.04 & 4 & 1 & 15 \\
    50$\times$15 & 20 & 25.18 & 22.81 & 17.87 & 17.64 & 0.23 & 3 & 1 & 14 \\
    40$\times$20 & 10 & 37.00 & 32.61 & 24.05 & 24.05 & 0.00 & 0 & 0 & 10 \\
    50$\times$20 & 20 & 29.27 & 27.05 & 21.99 & 21.94 & 0.05 & 3 & 1 & 15 \\
    100$\times$20 & 10 & 11.26 & 11.26 & 11.02 & 11.02 & 0.00 & 0 & 0 & 1 \\
    \bottomrule
  \end{tabular}
  \par\vspace{2pt}{\footnotesize C: Chuffed; G: Gecode; C2: Chuffed V2. Margen: gap de la mejor configuración individual del grupo menos gap del VBS.}
\end{table}


En las instancias de menos de 300 operaciones, Chuffed es mejor o igual que
las demás en casi todos los casos (Chuffed V2 gana solo en una), porque prueba
la optimalidad de la mayoría de las pequeñas y su búsqueda libre supera a la
estrategia de Chuffed V2. En $20 \times 15$ y $20 \times 20$ las victorias se
reparten, y a partir de $30 \times 15$ Chuffed V2 gana en la mayoría de las
instancias, con diferencias de hasta 16 puntos de \textit{gap} medio (en
$50 \times 10$). Gecode solo gana en tres instancias grandes de la familia
\texttt{ta}.

\begin{figure}[H]
  \centering
  \includegraphics{dispersion_configuraciones.pdf}
  \caption{\textit{Gap} de Chuffed frente al de Chuffed V2 en cada instancia,
  coloreado según la configuración que obtiene la mejor solución. Los puntos
  bajo la diagonal son instancias en que Chuffed V2 es mejor.}
  \label{fig:dispersion}
\end{figure}

Un comportamiento análogo aparece al agrupar las instancias por familia en la
Tabla~\ref{tab:por_familia}, donde Chuffed domina en las familias de instancias
pequeñas y medianas (\texttt{ft}, \texttt{la}, \texttt{abz}, \texttt{orb}), las
dos configuraciones de Chuffed se reparten las victorias en \texttt{swv},
\texttt{yn} y \texttt{ta}, y Chuffed V2 domina en \texttt{dmu}, con 65 de 80
victorias.

\begin{table}[H]
  \centering
  \small
  \singlespacing
  \caption{Resultados por familia de instancias.}
  \label{tab:por_familia}
  \begin{tabular}{lrrrrrrrrrrrrrr}
    \toprule
    Familia & Inst. & \multicolumn{3}{c}{Óptimos} & \multicolumn{4}{c}{Gap medio (\%)} & \multicolumn{3}{c}{Gana} & \multicolumn{3}{c}{Sin solución} \\
    \cmidrule(lr){3-5}\cmidrule(lr){6-9}\cmidrule(lr){10-12}\cmidrule(lr){13-15}
     &  & C & G & C2 & C & G & C2 & VBS & C & G & C2 & C & G & C2 \\
    \midrule
    ft & 3 & 2 & 1 & 1 & 0.80 & 12.72 & 7.77 & 0.80 & 2 & 0 & 0 & 0 & 0 & 0 \\
    la & 36 & 23 & 5 & 13 & 1.11 & 9.77 & 4.15 & 1.10 & 20 & 0 & 1 & 0 & 4 & 2 \\
    abz & 5 & 2 & 0 & 1 & 3.81 & 15.97 & 7.89 & 3.72 & 3 & 0 & 1 & 0 & 0 & 0 \\
    orb & 10 & 9 & 0 & 3 & 0.13 & 9.19 & 5.08 & 0.13 & 7 & 0 & 0 & 0 & 0 & 0 \\
    swv & 18 & 0 & 0 & 0 & 21.08 & 27.17 & 19.51 & 13.80 & 9 & 0 & 8 & 0 & 3 & 2 \\
    yn & 4 & 0 & 0 & 0 & 8.23 & 20.87 & 9.46 & 7.34 & 2 & 0 & 2 & 0 & 0 & 0 \\
    ta & 80 & 0 & 0 & 0 & 13.79 & 20.54 & 15.01 & 12.81 & 43 & 3 & 23 & 11 & 36 & 16 \\
    dmu & 80 & 0 & 0 & 0 & 29.70 & 29.25 & 21.32 & 20.81 & 14 & 0 & 65 & 1 & 22 & 4 \\
    \bottomrule
  \end{tabular}
  \par\vspace{2pt}{\footnotesize C: Chuffed; G: Gecode; C2: Chuffed V2.}
\end{table}


A partir de esta distribución se observa que acertar en la elección entre
Chuffed y Chuffed V2 basta para capturar casi todo el margen disponible, dado
que Gecode aporta muy poco al VBS. Sin embargo, la frontera entre ambas
configuraciones carece de un corte único, puesto que en los tamaños de
$20 \times 15$ a $30 \times 15$ las victorias están repartidas y las
características de la instancia, adicionales al tamaño, pueden inclinar la
decisión.

\subsection{Comportamiento en el tiempo}

La evolución del \textit{gap} medio de cada configuración según el tiempo
transcurrido puede seguirse en la Tabla~\ref{tab:anytime} y en la
Figura~\ref{fig:anytime}.

\begin{table}[H]
  \centering
  \small
  \singlespacing
  \caption{Gap medio (\%) y óptimos probados según el tiempo transcurrido.}
  \label{tab:anytime}
  \begin{tabular}{rrrrrrrr}
    \toprule
    Tiempo (s) & \multicolumn{4}{c}{Gap medio (\%)} & \multicolumn{3}{c}{Óptimos probados} \\
    \cmidrule(lr){2-5}\cmidrule(lr){6-8}
     & C & G & C2 & VBS & C & G & C2 \\
    \midrule
    10 & 22.19 & 22.25 & 16.82 & 16.02 & 22 & 4 & 9 \\
    30 & 20.53 & 21.99 & 16.11 & 15.04 & 26 & 5 & 13 \\
    60 & 19.36 & 21.94 & 15.76 & 14.46 & 28 & 5 & 16 \\
    120 & 18.21 & 21.79 & 15.42 & 13.77 & 30 & 6 & 16 \\
    300 & 16.76 & 21.68 & 15.08 & 12.83 & 36 & 6 & 18 \\
    \bottomrule
  \end{tabular}
  \par\vspace{2pt}{\footnotesize C: Chuffed; G: Gecode; C2: Chuffed V2.}
\end{table}


\begin{figure}[H]
  \centering
  \includegraphics{anytime.pdf}
  \caption{\textit{Gap} medio de cada configuración y del VBS en función del
  tiempo (escala logarítmica).}
  \label{fig:anytime}
\end{figure}

Chuffed V2 encuentra buenas soluciones muy pronto, pues a los 10 segundos su
\textit{gap} medio ya es de 16.8\,\%, cercano a su valor final, mientras que
Chuffed necesita los 300 segundos para llegar a 16.8\,\%. Gecode casi no
mejora después de los primeros segundos. Las dos configuraciones de Chuffed,
en cambio, continúan mejorando hasta el final del presupuesto, lo que anticipa
una pérdida de calidad potencial para cualquier estrategia que fraccione el
tiempo entre ellas.

%------------------------------------------------------------------------------
\section{\textit{Winner Takes All} frente al \textit{schedule} secuencial}
\label{sec:res_wta_schedule}
%------------------------------------------------------------------------------

La Tabla~\ref{tab:portafolios} y la Figura~\ref{fig:portafolios} comparan
todas las estrategias.

\subsection{Todas las estrategias superan al mejor \textit{solver} individual}

Todos los métodos de portafolio cierran entre el 76\,\% y el 87\,\% de la
brecha entre el SBS y el VBS, y la mejora frente al SBS es
estadísticamente significativa en todos los casos (por ejemplo, 1.96 puntos
para WTA $k$-NN, con un intervalo de confianza del 95\,\% de
$[1.48, 2.49]$, y 1.80 puntos para el \textit{schedule} con ML, con
$[1.36, 2.30]$). Asimismo, todos los métodos recuperan los 36 óptimos de
Chuffed o casi todos, que el SBS (Chuffed V2) pierde a la mitad. Incluso la
regla por tamaño, que usa una sola característica, cierra el 76\,\% de la
brecha.

\begin{table}[H]
  \centering
  \small
  \singlespacing
  \caption{Comparación de los \textit{solvers} sin cotas y con cotas y de los enfoques de portafolio (236 instancias, 300\,s por instancia, validación \textit{leave-one-out}).}
  \label{tab:portafolios}
  \begin{tabular}{lrrrrrr}
    \toprule
    Método & Gap (\%) & Óptimos & PAR10 (s) & Brecha cerrada (\%) & Mejor & Peor \\
    \midrule
    Chuffed (sin cotas) & 648.32 & 28 & 2652 & -- & 13 & 205 \\
    Gecode (sin cotas) & 25.57 & 5 & 2938 & -- & 7 & 225 \\
    VBS (sin cotas) & 24.53 & 28 & 2651 & -- & 17 & 201 \\
    \midrule
    Chuffed & 16.76 & 36 & 2550 & -- & 103 & 101 \\
    Gecode & 21.68 & 6 & 2924 & -- & 14 & 197 \\
    Chuffed V2 & 15.08 & 18 & 2774 & -- & 0 & 0 \\
    \midrule
    Regla por tamaño & 13.37 & 36 & 2550 & 76 & 68 & 8 \\
    WTA $k$-NN (mzn2feat) & 13.23 & 36 & 2550 & 82 & 84 & 18 \\
    WTA $k$-NN (propias) & 13.34 & 36 & 2550 & 77 & 83 & 17 \\
    WTA $k$-NN (ambas) & 13.12 & 36 & 2550 & 87 & 88 & 17 \\
    SUNNY (mzn2feat) & 13.20 & 36 & 2550 & 84 & 92 & 43 \\
    SUNNY (propias) & 13.28 & 34 & 2576 & 80 & 91 & 43 \\
    SUNNY (ambas) & 13.20 & 36 & 2552 & 84 & 92 & 40 \\
    \midrule
    \textit{Schedule} estático & 13.31 & 35 & 2567 & 79 & 97 & 70 \\
    \textit{Schedule} con ML (híbrido) & 13.28 & 35 & 2564 & 80 & 87 & 42 \\
    \midrule
    VBS & 12.83 & 36 & 2550 & 100 & 104 & 0 \\
    \bottomrule
  \end{tabular}
  \par\vspace{2pt}{\footnotesize Cada fila es un método. Gap: \textit{gap} medio frente a la mejor solución conocida. Óptimos: instancias en que se prueba la optimalidad. PAR10: tiempo medio hasta probar el óptimo, contando 3000\,s si no se prueba. En Gap y PAR10, un valor menor es mejor. Brecha cerrada: parte de la distancia entre el SBS (Chuffed V2) y el VBS que recupera el método, $100\cdot(\mathrm{gap}_{SBS}-\mathrm{gap})/(\mathrm{gap}_{SBS}-\mathrm{gap}_{VBS})$ (0 = SBS, 100 = VBS). Mejor/Peor: instancias con menor/mayor \textit{makespan} que el SBS. Sin cotas: Chuffed y Gecode con el modelo original y su VBS, es decir, el portafolio planteado al inicio del trabajo; Chuffed V2 no tiene fila sin cotas porque se definió como configuración del modelo con cotas. Las demás filas usan el modelo con cotas.}
\end{table}


Las filas sin cotas de la Tabla~\ref{tab:portafolios} corresponden al
portafolio planteado al inicio del trabajo, formado por Chuffed y Gecode con
el modelo original. Con ese modelo, el \textit{gap} medio de Chuffed era de
648.3\,\% y el de Gecode de 25.6\,\%, y aun eligiendo en cada instancia la
mejor de las dos, el VBS no bajaba del 24.5\,\%, casi el doble del 13\,\%
que alcanzan los enfoques de portafolio con el modelo con cotas. Además, el
margen entre el mejor \textit{solver} y el VBS era de apenas 1.0 puntos,
frente a 2.24 con cotas, de modo que el acotamiento del modelo es lo que le
da al portafolio espacio para mejorar. Chuffed V2 se incorporó después de
introducir las cotas, como una segunda estrategia de búsqueda sobre ese
modelo, por lo que no se evaluó con el modelo original.

\begin{figure}[H]
  \centering
  \includegraphics{portafolios_gap.pdf}
  \caption{\textit{Gap} medio de cada método. Las líneas punteadas marcan el
  SBS (gris) y el VBS (verde).}
  \label{fig:portafolios}
\end{figure}

\subsection{La selección supera en promedio al reparto del tiempo}

Con los datos de este trabajo, la mejor estrategia en promedio es
\textit{Winner Takes All} con las 85 características (13.12\,\%). El análisis
pareado frente a las demás estrategias, instancia por instancia
(Tabla~\ref{tab:comparaciones}), muestra que WTA obtiene un
\textit{makespan} igual o mejor que los \textit{schedules} en
la mayoría de las instancias en que difieren. La diferencia media frente a
SUNNY es pequeña y su intervalo de confianza incluye el cero, pero WTA es
mejor en 70 instancias y peor en 19. Frente al \textit{schedule} estático y
al \textit{schedule} con ML la ventaja es significativa, aunque de pocas
décimas.

\begin{table}[H]
  \centering
  \small
  \singlespacing
  \caption{Comparación pareada de WTA $k$-NN (ambas) con otros métodos en las 236 instancias.}
  \label{tab:comparaciones}
  \begin{tabular}{lrcr}
    \toprule
    Frente a & Diferencia & IC 95\,\% & WTA mejor / peor \\
    \midrule
    SUNNY (ambas) & $-0.08$ & $[-0.19, +0.04]$ & 70 / 19 \\
    \textit{Schedule} con ML (híbrido) & $-0.16$ & $[-0.28, -0.02]$ & 95 / 22 \\
    \textit{Schedule} estático & $-0.19$ & $[-0.33, -0.04]$ & 126 / 28 \\
    Regla por tamaño & $-0.25$ & $[-0.49, -0.03]$ & 26 / 15 \\
    \bottomrule
  \end{tabular}
  \par\vspace{2pt}{\footnotesize Diferencia: \textit{gap} medio de WTA menos el del otro método, en puntos porcentuales; un valor negativo favorece a WTA. IC 95\,\%: intervalo de confianza de esa diferencia, obtenido por \textit{bootstrap} con 10\,000 remuestreos de las instancias; si no incluye el cero, la diferencia es significativa. WTA mejor / peor: instancias en que WTA obtiene un \textit{makespan} menor / mayor que el otro método.}
\end{table}


Este comportamiento responde a dos factores observados en las secciones
previas. Primero, la selección sigue un patrón por tamaño muy marcado, lo que
permite a WTA obtener el mismo resultado que el VBS en 202 de las 236
instancias. Segundo, las dos configuraciones de Chuffed siguen mejorando hasta
los 300 segundos, de modo que reducir el tiempo asignado a la configuración
adecuada para destinarlo a otra disminuye la calidad alcanzada. Un
\textit{schedule} asume esa reducción en muchas instancias para protegerse de
un error de selección que, en este conjunto, ocurre pocas veces.

\subsection{El \textit{schedule} es más robusto}

El promedio constituye solo una dimensión del desempeño. La motivación de los
\textit{schedules} es reducir el riesgo de una mala selección, y la
Tabla~\ref{tab:robustez} muestra que en ese aspecto sí son mejores. Para cada
método se mide la pérdida frente al VBS, es decir, cuántos puntos de
\textit{gap} pierde en cada instancia respecto a la mejor configuración.

\begin{table}[H]
  \centering
  \small
  \singlespacing
  \caption{Pérdida frente al VBS (puntos de \textit{gap}) en las 236 instancias.}
  \label{tab:robustez}
  \begin{tabular}{lrrrrr}
    \toprule
    Método & Media & Percentil 90 & Máxima & $> 2$ puntos & $> 5$ puntos \\
    \midrule
    SBS (Chuffed V2) & 2.24 & 7.53 & 22.75 & 76 & 42 \\
    Regla por tamaño & 0.54 & 1.59 & 9.89 & 23 & 11 \\
    WTA $k$-NN (ambas) & 0.28 & 0.45 & 7.82 & 13 & 4 \\
    SUNNY (ambas) & 0.36 & 1.07 & 6.04 & 9 & 3 \\
    \textit{Schedule} con ML (híbrido) & 0.44 & 1.40 & 5.53 & 14 & 3 \\
    \textit{Schedule} estático & 0.47 & 1.28 & 4.31 & 9 & 0 \\
    \bottomrule
  \end{tabular}
\end{table}


WTA tiene la menor pérdida media, pero también la mayor pérdida máxima entre
los portafolios (7.82 puntos) y cuatro instancias en las que pierde más de 5
puntos, correspondientes a los casos en que elige la configuración equivocada.
Los \textit{schedules} cometen más errores pequeños, porque siempre ceden algo
de tiempo, pero limitan los grandes. El \textit{schedule} estático nunca pierde
más de 4.31 puntos, y SUNNY combina una pérdida media cercana a la de WTA con
menos instancias por encima de 2 puntos.

Respecto a la pregunta P3, la comparación refleja un compromiso entre el
desempeño promedio y el control del riesgo. Mientras que seleccionar una sola
configuración con WTA es ligeramente mejor en promedio, en el peor caso los
\textit{schedules} resultan más robustos, por lo cual la elección entre ambos
esquemas depende de si se prioriza el desempeño medio o acotar el efecto de una
mala decisión.

\subsection{Sensibilidad al número de vecinos}

Como se observa en la Tabla~\ref{tab:sensibilidad_k}, ni WTA ni SUNNY dependen
mucho de $k$, puesto que entre $k = 3$ y $k = 30$ el \textit{gap} medio varía
menos de 0.15 puntos. El valor $k = 15$ no se ajustó con estos resultados,
sino que se fijó de antemano como la raíz cuadrada del número de instancias
de entrenamiento, criterio que usa SUNNY-CP \citep{amadini2015sunnycp}, y
está entre los mejores para ambos métodos.

\begin{table}[H]
  \centering
  \small
  \singlespacing
  \caption{Sensibilidad al número de vecinos $k$ (características: ambas).}
  \label{tab:sensibilidad_k}
  \begin{tabular}{rrr}
    \toprule
    $k$ & Gap WTA (\%) & Gap SUNNY (\%) \\
    \midrule
    3 & 13.26 & 13.23 \\
    5 & 13.16 & 13.21 \\
    7 & 13.22 & 13.22 \\
    10 & 13.18 & 13.23 \\
    15 & 13.12 & 13.20 \\
    20 & 13.14 & 13.18 \\
    30 & 13.12 & 13.18 \\
    \bottomrule
  \end{tabular}
\end{table}


%------------------------------------------------------------------------------
\section{Aporte de las características y del aprendizaje automático}
\label{sec:res_prediccion}
%------------------------------------------------------------------------------

\subsection{Conjuntos de características}

En WTA, la unión de las características de \texttt{mzn2feat} y las propias
da el mejor resultado (13.12\,\%), seguida de \texttt{mzn2feat} sola
(13.23\,\%) y de las propias solas (13.34\,\%). Las diferencias son pequeñas y
afectan a pocas instancias, ya que frente a cada conjunto por separado el
resultado de la unión cambia en solo 11 instancias. En SUNNY, \texttt{mzn2feat}
y la unión empatan (13.20\,\%). Que las características genéricas y las propias
den resultados tan parecidos es coherente con la estructura de la
Sección~\ref{sec:res_solvers}, en la cual gran parte de la decisión depende del
tamaño, magnitud que ambos conjuntos describen bien.

Frente a la regla por tamaño, que usa solo el número de operaciones, WTA con
todas las características reduce el \textit{gap} en 0.25 puntos
(Tabla~\ref{tab:comparaciones}). Las características adicionales aportan,
sobre todo, en la zona de transición entre Chuffed y Chuffed V2.

\subsection{Calidad de los modelos de predicción}

En la Tabla~\ref{tab:ml_calidad} se detallan las métricas predictivas de los
tres modelos del \textit{schedule} con ML, medidas con sus predicciones fuera
de muestra.

\begin{table}[H]
  \centering
  \small
  \singlespacing
  \caption{Calidad de los modelos de aprendizaje automático (predicciones fuera de muestra, \textit{leave-one-out}).}
  \label{tab:ml_calidad}
  \begin{tabular}{llr}
    \toprule
    Modelo & Métrica & Valor \\
    \midrule
    Resolución (Chuffed) & AUC & 0.98 \\
    Resolución (Chuffed) & Exactitud (\%) & 95.3 \\
    Resolución (Gecode) & AUC & 0.98 \\
    Resolución (Gecode) & Exactitud (\%) & 97.0 \\
    Resolución (Chuffed V2) & AUC & 0.97 \\
    Resolución (Chuffed V2) & Exactitud (\%) & 94.1 \\
    Tiempo condicional (Chuffed) & Error absoluto medio (s) & 41.7 \\
    Tiempo condicional (Chuffed) & Error relativo mediano (\%) & 89.8 \\
    Mejor configuración & Exactitud (\%) & 76.7 \\
    Mejor configuración (Chuffed vs resto) & AUC & 0.85 \\
    Mejor configuración (Gecode vs resto) & AUC & 0.70 \\
    Mejor configuración (Chuffed V2 vs resto) & AUC & 0.86 \\
    \bottomrule
  \end{tabular}
\end{table}


El clasificador de resolución es muy preciso, con un área bajo la curva ROC
(AUC) de 0.97 a 0.98 y una exactitud de entre el 94\,\% y el 97\,\% para las
tres configuraciones, lo que indica que en este conjunto resulta sencillo
distinguir las instancias cuyos óptimos pueden probarse en 300 segundos. El
regresor de tiempo es el modelo más débil, con un error relativo mediano del
89.8\,\%, del mismo orden que el propio tiempo, razón por la cual el
\textit{schedule} asigna el doble del tiempo predicho. Su bajo desempeño
tiene varias causas posibles. Solo 60 de las 708 ejecuciones prueban el
óptimo y sirven para entrenarlo, casi todas en instancias pequeñas (100
operaciones de mediana, frente a 425 en el conjunto completo), y sus tiempos
van de 0.3 a 273 segundos. Además, el tiempo que tarda un \textit{solver} en
probar el óptimo depende en buena parte de su recorrido de búsqueda, que las
características estáticas de la instancia no describen, y las ejecuciones
que agotan los 300 segundos no indican cuánto habrían tardado. El clasificador de la
mejor configuración acierta en el 76.7\,\% de las instancias, con un AUC de
0.85 a 0.86 para las dos configuraciones de Chuffed y de 0.70 para Gecode, que
gana muy pocas veces.

\subsection{\textit{Schedule} con ML frente al \textit{schedule} estático}

La Tabla~\ref{tab:ml_vs_estatico} compara el \textit{schedule} con ML con el
\textit{schedule} estático, que reparte el tiempo igual en todas las
instancias.

\begin{table}[H]
  \centering
  \small
  \singlespacing
  \caption{\textit{Schedule} con ML frente al \textit{schedule} estático (diferencia ML $-$ estático).}
  \label{tab:ml_vs_estatico}
  \begin{tabular}{lr}
    \toprule
     & Valor \\
    \midrule
    Instancias con resultado distinto & 102 \\
    ML mejor / peor & 68 / 34 \\
    Diferencia media de gap (puntos) & -0.032 \\
    Intervalo de confianza del 95\,\% (\textit{bootstrap}) & [-0.145, +0.085] \\
    Mayor ganancia / mayor pérdida (puntos) & 4.24 / 4.98 \\
    Prueba de Wilcoxon ($p$) & 0.052 \\
    \bottomrule
  \end{tabular}
\end{table}


El \textit{schedule} con ML adapta el reparto a cada instancia (produce 181
\textit{schedules} distintos) y es mejor que el estático en 68 instancias y
peor en 34. Sin embargo, la diferencia media es de solo 0.03 puntos, su
intervalo de confianza incluye el cero y la prueba de Wilcoxon no es
significativa al 5\,\%. Los dos métodos son, por tanto, equivalentes en
promedio.

El \textit{schedule} con ML queda por detrás de WTA porque reparte el tiempo en
proporción a la probabilidad de que cada configuración sea la mejor, y esa
probabilidad rara vez es cercana a 1, de modo que incluso cuando el
clasificador acierta, cede parte del tiempo a las otras configuraciones. Se
probó también una versión del clasificador que pondera cada instancia de
entrenamiento por el costo de equivocarse, sin mejoras apreciables (13.26\,\%),
por lo que no se adoptó.

En relación con la pregunta P4, las características de la instancia aportan con
claridad, dado que todas las estrategias que las usan superan al SBS y a la
regla por tamaño. Aunque el aprendizaje automático produce modelos precisos, su
empleo para repartir el tiempo obtiene un desempeño medio equivalente al de un
reparto fijo bien elegido e inferior al de la selección por vecinos más
cercanos.

%------------------------------------------------------------------------------
\section{Comparación con el \textit{solver} de referencia}
\label{sec:res_referencia}
%------------------------------------------------------------------------------

CP-SAT se ejecutó con el modelo con cotas bajo el mismo protocolo que las
configuraciones del portafolio. Sus cifras globales ya se incluyeron como
referencia en la Tabla~\ref{tab:individual} (un \textit{gap} medio de
3.97\,\%, 62 óptimos probados y ninguna instancia sin solución), mientras que
el desglose frente al portafolio por número de operaciones se presenta en la
Tabla~\ref{tab:cpsat}.

\begin{table}[H]
  \centering
  \small
  \singlespacing
  \caption{CP-SAT frente al portafolio por número de operaciones (modelo con cotas, 300\,s).}
  \label{tab:cpsat}
  \begin{tabular}{lrrrrrrrr}
    \toprule
    Operaciones & Inst. & \multicolumn{4}{c}{Gap medio (\%)} & \multicolumn{2}{c}{Óptimos} & VBS mejor \\
    \cmidrule(lr){3-6}\cmidrule(lr){7-8}
     &  & CP-SAT & SBS & WTA $k$-NN & VBS & CP-SAT & VBS &  \\
    \midrule
    Hasta 150 & 36 & 0.00 & 3.25 & 0.13 & 0.13 & 34 & 32 & 0 \\
    151 a 300 & 58 & 1.49 & 13.28 & 8.20 & 7.88 & 18 & 4 & 0 \\
    301 a 600 & 82 & 4.13 & 18.58 & 17.91 & 17.37 & 7 & 0 & 0 \\
    Más de 600 & 60 & 8.52 & 19.13 & 19.13 & 19.04 & 3 & 0 & 1 \\
    \midrule
    Total & 236 & 3.97 & 15.08 & 13.12 & 12.83 & 62 & 36 & 1 \\
    \bottomrule
  \end{tabular}
  \par\vspace{2pt}{\footnotesize SBS: Chuffed V2. VBS mejor: instancias en que la mejor configuración del portafolio obtiene un \textit{makespan} menor que CP-SAT.}
\end{table}


CP-SAT supera ampliamente al portafolio. Su \textit{gap} medio es la tercera
parte del VBS (3.97\,\% frente a 12.83\,\%), prueba la optimalidad en 62
instancias frente a 36 y obtiene un \textit{makespan} igual o mejor que la
mejor configuración del portafolio en 235 de las 236 instancias, con la única
excepción de \texttt{dmu76}. La diferencia crece con el tamaño, ya que en las
instancias de hasta 150 operaciones todos los métodos están cerca del óptimo,
pero en las de más de 300 operaciones CP-SAT mantiene un \textit{gap} de 4 a
9\,\%, mientras que el portafolio se sitúa entre el 17 y el 19\,\%. Asimismo,
CP-SAT alcanza a los 10 segundos un \textit{gap} medio de 12.9\,\%, similar al
del portafolio a los 300 segundos.

Este resultado confirma lo observado con el modelo original
(Sección~\ref{sec:ejecucion_solvers}) y respalda la decisión de mantener a
CP-SAT fuera del portafolio, puesto que el VBS calculado junto con CP-SAT
coincide con este último en todas las instancias salvo una. La ventaja de
CP-SAT proviene de su arquitectura, que combina búsqueda con aprendizaje de
cláusulas, relajaciones lineales y búsqueda local sobre las soluciones
encontradas \citep{perron2023}. La comparación muestra, por tanto, que la
selección entre configuraciones de \textit{solvers} complementarios mejora de
forma consistente al mejor de ellos, aunque sin compensar la diferencia frente
a un \textit{solver} de nivel superior.

%------------------------------------------------------------------------------
\section{Discusión}
\label{sec:discusion}
%------------------------------------------------------------------------------

\subsection{Respuesta a las preguntas de evaluación}

\begin{enumerate}[label=\textbf{P\arabic*.}, leftmargin=*]
  \item \textbf{Efecto de las cotas.} Es el factor que más influye en el
    desempeño. Reduce el \textit{gap} de Chuffed de 648\,\% a 17\,\% y lo
    convierte en la base del portafolio, mientras que en Gecode la mejora es
    moderada.
  \item \textbf{Complementariedad.} Las configuraciones son complementarias y
    muestran un patrón bien definido, con ventaja de Chuffed en las instancias
    pequeñas y medianas, de Chuffed V2 en las grandes y un aporte marginal de
    Gecode. El margen entre el SBS y el VBS es de 2.24 puntos.
  \item \textbf{WTA frente a \textit{schedule}.} En promedio, con estas
    configuraciones, la selección de una configuración con WTA es la mejor estrategia (87\,\% de la brecha
    cerrada), en tanto que los \textit{schedules} quedan algo por detrás en el
    promedio pero resultan más robustos frente a errores de selección.
  \item \textbf{Características y aprendizaje automático.} Las características
    permiten cerrar entre el 77\,\% y el 87\,\% de la brecha. Los modelos de
    aprendizaje son precisos, pero el \textit{schedule} construido con ellos
    es equivalente al estático.
  \item \textbf{Estado del arte.} CP-SAT, usado como referencia externa,
    supera a este portafolio en casi todas las instancias (3.97\,\% de
    \textit{gap} frente a 13.12\,\%). El portafolio mejora a sus miembros,
    pero no alcanza a CP-SAT.
\end{enumerate}

\subsection{Amenazas a la validez}

\begin{itemize}
  \item El margen disponible entre el SBS y el VBS es de 2.24 puntos de
    \textit{gap}, y las diferencias entre los métodos de portafolio son de
    décimas, razón por la cual las comparaciones se apoyan en pruebas pareadas
    y no solo en promedios.
  \item Todos los resultados se obtienen sobre instancias clásicas de la
    literatura (236 en la evaluación del portafolio) y, aunque la validación \textit{leave-one-out} evita evaluar
    con instancias de entrenamiento, las conclusiones podrían variar en
    instancias de otra distribución o estructura.
  \item Los candidatos del \textit{schedule} estático y la regla de reparto del
    \textit{schedule} con ML se definieron tras analizar la campaña, por lo que
    aunque la elección final se calcula dentro de cada partición de
    entrenamiento, el conjunto de alternativas estuvo guiado por los datos.
  \item Las mejores soluciones conocidas (BKS) corresponden a los resultados
    recopilados hasta 2020 \citep{weise2019jssp}, por lo cual en las instancias
    abiertas el \textit{gap} es una cota inferior de la distancia real al
    óptimo, condición que afecta por igual a todos los métodos.
  \item Las configuraciones dentro de un \textit{schedule} se ejecutan sin
    compartir información, de modo que un portafolio que transmitiera la mejor
    solución de una configuración a la siguiente como nueva cota superior
    podría obtener mejores resultados que los reportados.
  \item Cada configuración se ejecutó una sola vez por instancia y, aunque los
    \textit{solvers} son deterministas con parámetros fijos, los tiempos de
    pared presentan pequeñas fluctuaciones que pueden alterar la última
    solución alcanzada cerca del límite.
\end{itemize}

%==============================================================================
\chapter{Conclusiones y trabajo futuro}
\label{cap:conclusiones}
%==============================================================================

%------------------------------------------------------------------------------
\section{Conclusiones}
\label{sec:conclusiones}
%------------------------------------------------------------------------------

A lo largo de este trabajo se desarrolló un \textit{schedule} de ejecución
basado en portafolios para el JSSP clásico, que decide para cada instancia qué
configuraciones ejecutar, en qué orden y durante cuánto tiempo, a partir de
las características de la instancia y de modelos de aprendizaje automático.
Las conclusiones se organizan según los objetivos específicos y se refieren
al portafolio estudiado, formado por tres configuraciones de Chuffed y
Gecode, a las 236 instancias evaluadas y a un presupuesto de 300 segundos.

\paragraph{Objetivo 1: comparar enfoques de portafolio.} Se implementaron y
compararon, con la misma base de conocimiento y el mismo esquema de
validación, la selección de una configuración (\textit{Winner Takes All} con
vecinos más cercanos), el \textit{schedule} secuencial de SUNNY adaptado a
optimización, un \textit{schedule} estático y un \textit{schedule} construido
con aprendizaje automático. En el conjunto estudiado, seleccionar una sola
configuración da el mejor resultado en promedio, con un \textit{gap} medio de
13.12\,\% y el 87\,\% de la brecha entre el mejor \textit{solver} individual y
el VBS cerrada. Los \textit{schedules} quedan unas décimas por detrás en
promedio, pero son más robustos, puesto que el \textit{schedule} estático nunca
pierde más de 4.3 puntos de \textit{gap} frente a la mejor configuración,
mientras que \textit{Winner Takes All} llega a perder 7.8 cuando se equivoca. La
elección entre ambos enfoques depende, por tanto, de si se prioriza el
desempeño medio o la protección frente a errores de selección. La ventaja de
la selección se explica porque, en este conjunto, la mejor configuración es
fácil de predecir y las configuraciones siguen mejorando hasta el final del
presupuesto, con lo cual fraccionar el tiempo entre varias ejecuciones reduce el
margen de mejora de la configuración principal.

\paragraph{Objetivo 2: características de las instancias.} Se construyó un
conjunto de datos con 65 características genéricas de \texttt{mzn2feat} y 20
características propias del JSSP, sobre la carga de las máquinas, las
duraciones y la estructura de las rutas. Las estrategias que usan
características cierran entre el 77\,\% y el 87\,\% de la brecha, frente al
76\,\% de una regla que solo usa el tamaño de la instancia. Los dos grupos de
características dan resultados parecidos y su unión es la mejor opción para
\textit{Winner Takes All}. Las características propias, que se calculan en
menos de un segundo sin herramientas externas, constituyen una alternativa
práctica a \texttt{mzn2feat}.

\paragraph{Objetivo 3: portafolio con predicción y selección automática.} Se
implementó un prototipo funcional que, para una instancia nueva, calcula sus
cotas, extrae sus características, estima con tres modelos de bosques
aleatorios si cada configuración probará el óptimo, en cuánto tiempo y cuál
obtendrá la mejor solución, y construye con esas predicciones el orden y el
tiempo de ejecución. Los modelos son precisos (un AUC de 0.97 a 0.98 para
predecir la resolución y de 0.85 a 0.86 para identificar cuándo gana cada
una de las configuraciones de Chuffed), pero el \textit{schedule} construido con ellos resultó
equivalente en promedio al \textit{schedule} estático. En consecuencia, con
este portafolio, predecir bien el desempeño individual no se tradujo de forma
directa en un mejor reparto del tiempo.

\paragraph{Objetivo 4: validación experimental.} El portafolio se validó con
236 instancias clásicas de ocho familias, mediante validación
\textit{leave-one-out} y la simulación exacta de los \textit{schedules} a
partir de las trayectorias de los \textit{solvers}. Todas las estrategias de
portafolio superan al mejor \textit{solver} individual por entre 1.7 y 2.0
puntos de \textit{gap}, una diferencia estadísticamente significativa, y
recuperan todos o casi todos los óptimos que este pierde.
Como referencia externa, OR-Tools CP-SAT con el mismo modelo
obtiene un \textit{gap} medio de 3.97\,\%, muy por debajo de las estrategias
de este portafolio, lo que muestra que, en este caso, la selección entre estas
configuraciones mejora de forma consistente a sus miembros, pero no compensa
la distancia frente a CP-SAT.

\paragraph{Importancia de la formulación del modelo.} La mejora de mayor
magnitud en los experimentos provino de la formulación acotada del problema en
MiniZinc. Acotar el \textit{makespan} con una cota inferior y con la solución de
una heurística de despacho redujo el \textit{gap} medio de Chuffed de 648\,\% a
17\,\%, una ganancia muy superior a la que aporta cualquier estrategia de
selección sobre las mismas configuraciones. La comparación con los mismos
\textit{solvers} sin cotas (Tabla~\ref{tab:portafolios}) lo confirma: sin
cotas, ni eligiendo en cada instancia la mejor entre Chuffed y Gecode el
\textit{gap} medio baja del 24.5\,\%, y el margen entre el mejor
\textit{solver} y el VBS es de apenas 1.0 punto, mientras que con cotas los
enfoques de portafolio quedan cerca del 13\,\% con un margen de 2.24 puntos.
Asimismo, la heurística proporciona
una solución de respaldo que resguarda al portafolio frente a configuraciones
que no encuentran solución, y en seis instancias prueba la optimalidad sin
ejecutar ningún \textit{solver}. En este portafolio, aprovechar la
complementariedad entre sus miembros dependió de contar antes con una
formulación bien acotada para cada uno de ellos.

%------------------------------------------------------------------------------
\section{Limitaciones}
\label{sec:limitaciones}
%------------------------------------------------------------------------------

\begin{itemize}
  \item El portafolio activo se compone únicamente de dos \textit{solvers} de
    programación por restricciones en tres configuraciones y, aunque la
    heurística de despacho participa mediante el cálculo de las cotas y como
    solución de respaldo, no se incluyeron las metaheurísticas contempladas en
    el anteproyecto dentro del conjunto de algoritmos seleccionables.
  \item La evaluación se realizó sobre 236 instancias clásicas de la
    literatura, de modo que las conclusiones acerca de qué configuración
    conviene en cada caso están ligadas a esa distribución particular de tamaños
    y estructuras.
  \item El margen disponible entre el mejor \textit{solver} individual y el VBS
    es de 2.24 puntos de \textit{gap}, lo que acota la ganancia máxima que
    puede alcanzar cualquier estrategia de selección con estas configuraciones.
  \item Todas las conclusiones corresponden a un presupuesto fijo de 300
    segundos por instancia, por lo cual el reparto óptimo del tiempo y la
    conveniencia relativa de un \textit{schedule} podrían variar bajo otros
    límites de cómputo.
  \item Las configuraciones que integran un \textit{schedule} se ejecutan de
    manera independiente, sin transmitir la mejor solución encontrada de una
    etapa a la siguiente.
  \item CP-SAT se mantuvo fuera del portafolio como referencia externa con el
    propósito de estudiar la selección entre configuraciones de desempeño
    complementario.
\end{itemize}

%------------------------------------------------------------------------------
\section{Trabajo futuro}
\label{sec:trabajo_futuro}
%------------------------------------------------------------------------------

\begin{itemize}
  \item Transmitir la mejor solución encontrada por una configuración a la
    siguiente como nueva cota superior o como solución inicial, de manera que
    las etapas sucesivas del \textit{schedule} aprovechen el avance previo y el
    reparto del tiempo pueda superar a la selección de una sola configuración.
  \item Ampliar la composición del portafolio incorporando metaheurísticas como
    la búsqueda tabú, cuyo comportamiento en instancias grandes difiere del de
    los \textit{solvers} de programación por restricciones, junto con nuevas
    estrategias de búsqueda para los \textit{solvers} actuales.
  \item Emplear generadores de instancias que cubran regiones del espacio de
    características poco representadas en el conjunto clásico, especialmente
    aquellas donde la mejor configuración sea más difícil de predecir y un
    \textit{schedule} secuencial tenga mayor margen de aporte.
  \item Diseñar reglas de asignación de tiempo basadas en la predicción directa
    del \textit{gap} esperado en función del tiempo, en lugar de repartir el
    presupuesto en proporción a la probabilidad de que cada configuración sea la
    mejor.
  \item Extraer los descriptores de \texttt{mzn2feat} directamente sobre el
    modelo con cotas, de modo que las características reflejen el problema
    acotado que resuelven los \textit{solvers}.
  \item Extender el esquema de portafolio a variantes como el \textit{Flexible
    Job Shop} y a formulaciones multiobjetivo del problema, que quedaron fuera
    del alcance de este trabajo.
\end{itemize}


%------------------------------------------------------------------------------
%  Referencias
%------------------------------------------------------------------------------
%==============================================================================
%  Referencias (formato APA, 7.a edición)
%  Cada \bibitem lleva entre corchetes la etiqueta que aparece en el texto:
%    \bibitem[Autor(Año)]{clave}           -> (Autor, Año)
%    \bibitem[{Autor1 \& Autor2}(Año)]{..}  -> (Autor1 & Autor2, Año)
%    \bibitem[{Autor1 et~al.}(Año)]{..}     -> (Autor1 et al., Año)
%  Mantener la lista en orden alfabético.
%==============================================================================
\cleardoublepage
\phantomsection
\addcontentsline{toc}{chapter}{\bibname}

\begin{thebibliography}{99}

\bibitem[{Adams et~al.}(1988)]{adams1988}
Adams, J., Balas, E., \& Zawack, D. (1988). The shifting bottleneck procedure
for job shop scheduling. \textit{Management Science}, \textit{34}(3), 391--401.

\bibitem[{Amadini et~al.}(2014a)]{amadini2014mzn2feat}
Amadini, R., Gabbrielli, M., \& Mauro, J. (2014a). An enhanced features
extractor for a portfolio of constraint solvers. En \textit{Proceedings of the
29th Annual ACM Symposium on Applied Computing (SAC 2014)} (pp. 1357--1359).
ACM.

\bibitem[{Amadini et~al.}(2014b)]{amadini2014sunny}
Amadini, R., Gabbrielli, M., \& Mauro, J. (2014b). SUNNY: A lazy portfolio
approach for constraint solving. \textit{Theory and Practice of Logic
Programming}, \textit{14}(4--5), 509--524.
\url{https://doi.org/10.1017/S1471068414000179}

\bibitem[{Amadini et~al.}(2015)]{amadini2015sunnycp}
Amadini, R., Gabbrielli, M., \& Mauro, J. (2015). SUNNY-CP: A sequential CP
portfolio solver. En \textit{Proceedings of the 30th Annual ACM Symposium on
Applied Computing (SAC 2015)}. ACM. \url{https://doi.org/10.1145/2695664.2695741}

\bibitem[{Applegate \& Cook}(1991)]{applegate1991}
Applegate, D., \& Cook, W. (1991). A computational study of the job-shop
scheduling problem. \textit{ORSA Journal on Computing}, \textit{3}(2),
149--156.

\bibitem[{Beasley}(s.~f.)]{orlibjobshop}
Beasley, J. E. (s. f.). \textit{OR-Library: Job shop scheduling} [Conjunto de
datos, archivo jobshop1.txt]. Recuperado el 10 de noviembre de 2026, de
\url{https://people.brunel.ac.uk/~mastjjb/jeb/orlib/jobshopinfo.html}

\bibitem[Beasley(1990)]{beasley1990}
Beasley, J. E. (1990). OR-Library: Distributing test problems by electronic
mail. \textit{Journal of the Operational Research Society}, \textit{41}(11),
1069--1072.

\bibitem[{Bengio et~al.}(2021)]{bengio2021}
Bengio, Y., Lodi, A., \& Prouvost, A. (2021). Machine learning for
combinatorial optimization: A methodological tour d'horizon. \textit{European
Journal of Operational Research}, \textit{290}(2), 405--421.

\bibitem[{Bischl et~al.}(2016)]{bischl2016aslib}
Bischl, B., Kerschke, P., Kotthoff, L., Lindauer, M., Malitsky, Y., Fréchette,
A., Hoos, H., Hutter, F., Leyton-Brown, K., Tierney, K., \& Vanschoren, J.
(2016). ASlib: A benchmark library for algorithm selection.
\textit{Artificial Intelligence}, \textit{237}, 41--58.

\bibitem[Breiman(2001)]{breiman2001}
Breiman, L. (2001). Random forests. \textit{Machine Learning}, \textit{45}(1),
5--32.

\bibitem[{Brucker et~al.}(1994)]{brucker1994}
Brucker, P., Jurisch, B., \& Sievers, B. (1994). A branch and bound algorithm
for the job-shop scheduling problem. \textit{Discrete Applied Mathematics},
\textit{49}(1--3), 107--127.

\bibitem[{Carlier \& Pinson}(1989)]{carlier1989}
Carlier, J., \& Pinson, E. (1989). An algorithm for solving the job-shop
problem. \textit{Management Science}, \textit{35}(2), 164--176.

\bibitem[{Carlier \& Pinson}(1994)]{carlier1994}
Carlier, J., \& Pinson, E. (1994). Adjustment of heads and tails for the
job-shop problem. \textit{European Journal of Operational Research},
\textit{78}(2), 146--161.

\bibitem[{Carrera-Rivera et~al.}(2022)]{carrera2022}
Carrera-Rivera, A., Ochoa, W., Larrinaga, F., \& Lasa, G. (2022). How-to
conduct a systematic literature review: A quick guide for computer science
research. \textit{MethodsX}, \textit{9}, 101895.
\url{https://doi.org/10.1016/j.mex.2022.101895}

\bibitem[Chu(2011)]{chu2011}
Chu, G. (2011). \textit{Improving combinatorial optimization} [Tesis
doctoral]. University of Melbourne.

\bibitem[COSEAL(2024)]{coseal2024}
COSEAL. (2024). \textit{COSEAL: Configuration and Selection of Algorithms}.
\url{https://www.coseal.net/}

\bibitem[{Cover \& Hart}(1967)]{cover1967}
Cover, T., \& Hart, P. (1967). Nearest neighbor pattern classification.
\textit{IEEE Transactions on Information Theory}, \textit{13}(1), 21--27.

\bibitem[{Demirkol et~al.}(1998)]{demirkol1998}
Demirkol, E., Mehta, S., \& Uzsoy, R. (1998). Benchmarks for shop scheduling
problems. \textit{European Journal of Operational Research}, \textit{109}(1),
137--141.

\bibitem[{Fisher \& Thompson}(1963)]{fisher1963}
Fisher, H., \& Thompson, G. L. (1963). Probabilistic learning combinations of
local job-shop scheduling rules. En J. F. Muth \& G. L. Thompson (Eds.),
\textit{Industrial scheduling} (pp. 225--251). Prentice-Hall.

\bibitem[Friedman(2001)]{friedman2001}
Friedman, J. H. (2001). Greedy function approximation: A gradient boosting
machine. \textit{The Annals of Statistics}, \textit{29}(5), 1189--1232.

\bibitem[{Garey et~al.}(1976)]{garey1976}
Garey, M. R., Johnson, D. S., \& Sethi, R. (1976). The complexity of flowshop
and jobshop scheduling. \textit{Mathematics of Operations Research},
\textit{1}(2), 117--129.

\bibitem[{Gecode Team}(2019)]{gecode2019}
Gecode Team. (2019). \textit{Gecode: Generic constraint development
environment} (versión 6) [Software]. \url{https://www.gecode.org}

\bibitem[{Giffler \& Thompson}(1960)]{giffler1960}
Giffler, B., \& Thompson, G. L. (1960). Algorithms for solving
production-scheduling problems. \textit{Operations Research}, \textit{8}(4),
487--503.

\bibitem[{Gomes \& Selman}(2001)]{gomes2001}
Gomes, C. P., \& Selman, B. (2001). Algorithm portfolios. \textit{Artificial
Intelligence}, \textit{126}(1--2), 43--62.

\bibitem[Google(2025)]{perron2025}
Google. (2025). \textit{OR-Tools: CP-SAT solver} (versión 9.12) [Software].
\url{https://developers.google.com/optimization/cp/}

\bibitem[{Hoos et~al.}(2015)]{hoos2015aspeed}
Hoos, H., Kaminski, R., Lindauer, M., \& Schaub, T. (2015). aspeed: Solver
scheduling via answer set programming. \textit{Theory and Practice of Logic
Programming}, \textit{15}(1), 117--142.
\url{https://doi.org/10.1017/S1471068414000015}

\bibitem[{Hutter et~al.}(2014)]{hutter2014}
Hutter, F., Xu, L., Hoos, H. H., \& Leyton-Brown, K. (2014). Algorithm runtime
prediction: Methods \& evaluation. \textit{Artificial Intelligence},
\textit{206}, 79--111.

\bibitem[{Kerschke et~al.}(2019)]{kerschke2019}
Kerschke, P., Hoos, H. H., Neumann, F., \& Trautmann, H. (2019). Automated
algorithm selection: Survey and perspectives. \textit{Evolutionary
Computation}, \textit{27}(1), 3--45.

\bibitem[Kotthoff(2014)]{kotthoff2014}
Kotthoff, L. (2014). Algorithm selection for combinatorial search problems: A
survey. \textit{AI Magazine}, \textit{35}(3), 48--60.

\bibitem[{Ku \& Beck}(2016)]{ku2016}
Ku, W.-Y., \& Beck, J. C. (2016). Mixed integer programming models for job
shop scheduling: A computational analysis. \textit{Computers \& Operations
Research}, \textit{73}, 165--173.

\bibitem[Lawrence(1984)]{lawrence1984}
Lawrence, S. (1984). \textit{Resource constrained project scheduling: An
experimental investigation of heuristic scheduling techniques (Supplement)}.
Graduate School of Industrial Administration, Carnegie Mellon University.

\bibitem[{Lindauer et~al.}(2019)]{lindauer2019}
Lindauer, M., van Rijn, J. N., \& Kotthoff, L. (2019). The algorithm
selection competitions 2015 and 2017. \textit{Artificial Intelligence},
\textit{272}, 86--100. \url{https://doi.org/10.1016/j.artint.2018.10.004}

\bibitem[{Mirshekarian \& Šormaz}(2016)]{mirshekarian2016}
Mirshekarian, S., \& Šormaz, D. N. (2016). Correlation of job-shop scheduling
problem features with scheduling efficiency. \textit{Expert Systems with
Applications}, \textit{62}, 131--147.

\bibitem[{Nethercote et~al.}(2007)]{nethercote2007}
Nethercote, N., Stuckey, P. J., Becket, R., Brand, S., Duck, G. J., \& Tack,
G. (2007). MiniZinc: Towards a standard CP modelling language. En C. Bessière
(Ed.), \textit{Principles and Practice of Constraint Programming -- CP 2007}
(pp. 529--543). Springer.

\bibitem[{Nowicki \& Smutnicki}(1996)]{nowicki1996}
Nowicki, E., \& Smutnicki, C. (1996). A fast taboo search algorithm for the
job shop problem. \textit{Management Science}, \textit{42}(6), 797--813.

\bibitem[{Ohrimenko et~al.}(2009)]{ohrimenko2009}
Ohrimenko, O., Stuckey, P. J., \& Codish, M. (2009). Propagation via lazy
clause generation. \textit{Constraints}, \textit{14}(3), 357--391.

\bibitem[{O'Mahony et~al.}(2008)]{omahony2008cphydra}
O'Mahony, E., Hebrard, E., Holland, A., Nugent, C., \& O'Sullivan, B. (2008).
Using case-based reasoning in an algorithm portfolio for constraint solving.
En \textit{Proceedings of the 19th Irish Conference on Artificial Intelligence
and Cognitive Science (AICS 2008)}.

\bibitem[{Panwalkar \& Iskander}(1977)]{panwalkar1977}
Panwalkar, S. S., \& Iskander, W. (1977). A survey of scheduling rules.
\textit{Operations Research}, \textit{25}(1), 45--61.

\bibitem[{Pedregosa et~al.}(2011)]{pedregosa2011}
Pedregosa, F., Varoquaux, G., Gramfort, A., Michel, V., Thirion, B., Grisel,
O., Blondel, M., Prettenhofer, P., Weiss, R., Dubourg, V., Vanderplas, J.,
Passos, A., Cournapeau, D., Brucher, M., Perrot, M., \& Duchesnay, É. (2011).
Scikit-learn: Machine learning in Python. \textit{Journal of Machine Learning
Research}, \textit{12}, 2825--2830.

\bibitem[{Perron et~al.}(2023)]{perron2023}
Perron, L., Didier, F., \& Gay, S. (2023). The CP-SAT-LP solver (invited
talk). En R. H. C. Yap (Ed.), \textit{29th International Conference on
Principles and Practice of Constraint Programming (CP 2023)} (LIPIcs, Vol.
280, pp. 3:1--3:2). Schloss Dagstuhl -- Leibniz-Zentrum für Informatik.
\url{https://doi.org/10.4230/LIPIcs.CP.2023.3}

\bibitem[Pinedo(2022)]{pinedo2022}
Pinedo, M. L. (2022). \textit{Scheduling: Theory, algorithms, and systems}
(6.a ed.). Springer.

\bibitem[{Reijnen et~al.}(2023)]{reijnen2023}
Reijnen, R., Smit, I. G., Zhang, H., Wu, Y., Bukhsh, Z., \& Zhang, Y. (2023).
\textit{Job shop scheduling benchmark: Environments and instances for learning
and non-learning methods} (arXiv:2308.12794). arXiv.
\url{https://arxiv.org/abs/2308.12794}

\bibitem[Rice(1976)]{rice1976}
Rice, J. R. (1976). The algorithm selection problem. \textit{Advances in
Computers}, \textit{15}, 65--118.

\bibitem[{Rossi et~al.}(2006)]{rossi2006}
Rossi, F., van Beek, P., \& Walsh, T. (Eds.). (2006). \textit{Handbook of
constraint programming}. Elsevier.

\bibitem[{Schutt et~al.}(2011)]{schutt2011}
Schutt, A., Feydy, T., Stuckey, P. J., \& Wallace, M. G. (2011). Explaining
the cumulative propagator. \textit{Constraints}, \textit{16}(3), 250--282.

\bibitem[{Smit et~al.}(2025)]{smit2024gnn}
Smit, I. G., Zhou, J., Reijnen, R., Wu, Y., Chen, J., Zhang, C., Bukhsh, Z.,
Zhang, Y., \& Nuijten, W. (2025). Graph neural networks for job shop
scheduling problems: A survey. \textit{Computers \& Operations Research},
\textit{176}, 106914. \url{https://doi.org/10.1016/j.cor.2024.106914}

\bibitem[Smith-Miles(2008)]{smithmiles2009}
Smith-Miles, K. (2008). Cross-disciplinary perspectives on meta-learning for
algorithm selection. \textit{ACM Computing Surveys}, \textit{41}(1), Artículo 6.
\url{https://doi.org/10.1145/1456650.1456656}

\bibitem[{Storer et~al.}(1992)]{storer1992}
Storer, R. H., Wu, S. D., \& Vaccari, R. (1992). New search spaces for
sequencing problems with application to job shop scheduling.
\textit{Management Science}, \textit{38}(10), 1495--1509.

\bibitem[{Strassl \& Musliu}(2022)]{strassl2022}
Strassl, S., \& Musliu, N. (2022). Instance space analysis and algorithm
selection for the job shop scheduling problem. \textit{Computers \&
Operations Research}, \textit{141}, 105661.

\bibitem[{Streeter \& Smith}(2008)]{streeter2008}
Streeter, M., \& Smith, S. F. (2008). New techniques for algorithm portfolio
design. En \textit{Proceedings of the 24th Conference on Uncertainty in
Artificial Intelligence (UAI 2008)} (pp. 519--527). AUAI Press.

\bibitem[{Stuckey et~al.}(2014)]{stuckey2014}
Stuckey, P. J., Feydy, T., Schutt, A., Tack, G., \& Fischer, J. (2014). The
MiniZinc Challenge 2008--2013. \textit{AI Magazine}, \textit{35}(2), 55--60.

\bibitem[Taillard(1993)]{taillard1993}
Taillard, E. (1993). Benchmarks for basic scheduling problems.
\textit{European Journal of Operational Research}, \textit{64}(2), 278--285.

\bibitem[Weise(2019)]{weise2019jssp}
Weise, T. (2019). \textit{jsspInstancesAndResults: Results, data, and
instances of the job shop scheduling problem} [Repositorio de software].
GitHub. \url{https://github.com/thomasWeise/jsspInstancesAndResults}

\bibitem[{Wolpert \& Macready}(1997)]{wolpert1997}
Wolpert, D. H., \& Macready, W. G. (1997). No free lunch theorems for
optimization. \textit{IEEE Transactions on Evolutionary Computation},
\textit{1}(1), 67--82.

\bibitem[{Xu et~al.}(2010)]{xu2010hydra}
Xu, L., Hoos, H. H., \& Leyton-Brown, K. (2010). Hydra: Automatically
configuring algorithms for portfolio-based selection. En \textit{Proceedings of
the AAAI Conference on Artificial Intelligence}, \textit{24}(1), 210--216.

\bibitem[{Xu et~al.}(2008)]{xu2008satzilla}
Xu, L., Hutter, F., Hoos, H. H., \& Leyton-Brown, K. (2008). SATzilla:
Portfolio-based algorithm selection for SAT. \textit{Journal of Artificial
Intelligence Research}, \textit{32}, 565--606.

\bibitem[{Xu et~al.}(2025)]{xu2025learn}
Xu, M., Mei, Y., Zhang, F., \& Zhang, M. (2025). Learn to optimise for job
shop scheduling: A survey with comparison between genetic programming and
reinforcement learning. \textit{Artificial Intelligence Review}, \textit{58},
160. \url{https://doi.org/10.1007/s10462-024-11059-9}

\bibitem[{Yamada \& Nakano}(1992)]{yamada1992}
Yamada, T., \& Nakano, R. (1992). A genetic algorithm applicable to
large-scale job-shop problems. En R. Männer \& B. Manderick (Eds.),
\textit{Parallel Problem Solving from Nature 2} (pp. 281--290). Elsevier.

\end{thebibliography}


%------------------------------------------------------------------------------
%  Anexos
%------------------------------------------------------------------------------
\appendix
\input{anexos/A_resultados_instancia}
\input{anexos/B_caracteristicas}
\input{anexos/C_configuracion}
\input{anexos/D_reproduccion}

\end{document}
